\documentclass{book}

\usepackage[bookmarks,bookmarksopen,bookmarksdepth=3]{hyperref}
\usepackage{amsmath}
\usepackage{amssymb}
\usepackage{amsthm}
\usepackage[all,cmtip]{xy}

\hypersetup{
  colorlinks=true,
  urlcolor=blue,
  linkcolor=magenta,
  citecolor=blue,
  filecolor=blue
}

\theoremstyle{definition}
\renewcommand{\proofname}{Proof}
\newenvironment{Proof}[1][Proof]
{\proof[#1]\leftskip=1cm\rightskip=1cm}{\endproof}

\newtheorem*{defn}{Definition}
\newtheorem{defs}{Definitions}
\newtheorem*{lema}{Lemma}
\newtheorem*{obs}{Observation}
\newtheorem*{teo}{Theorem}
\newtheorem*{prop}{Proposition}
\newtheorem*{coro}{Corollary}
\newtheorem*{ejer}{Exercise}
\newtheorem*{ejem}{Example}
\newtheorem*{ejems}{Examples}
\newtheorem*{pregunta}{Question}

\newcommand{\R}{\mathbb{R}}
\newcommand{\Z}{\mathbb{Z}}
\newcommand{\N}{\mathbb{N}}
\newcommand{\C}{\mathbb{C}}
\newcommand{\Q}{\mathbb{Q}}
\newcommand{\T}{\mathbb{T}}
\newcommand{\K}{\mathbb{K}}
\DeclareMathOperator{\coker}{coker}
\DeclareMathOperator{\img}{img}
\DeclareMathOperator{\ran}{ran}
\DeclareMathOperator{\Top}{Top}
\DeclareMathOperator{\Grp}{Grp}
\DeclareMathOperator{\rel}{rel}
\DeclareMathOperator{\Orb}{Orb}
\DeclareMathOperator{\Int}{int}

\title{Algebraic Topology Notes}
\author{Prof. Luis Jorge S\'anchez Salda\~na\\
Notes by Dani\\
\href{https://github.com/danimalabares/top-alg}{github.com/danimalabares/top-alg}}

\begin{document}
\maketitle
\phantomsection
\addcontentsline{toc}{part}{Contents}
\tableofcontents

\part{Fundamental group}

\chapter{The fundamental group}

\section{Paths and homotopies}

We denote by $I$ the unit interval $[0,1]$ and by
$X$ an arbitrary topological space.

\begin{defn}
A \textbf{path} is a function $f:I\to X$.
\end{defn}

\begin{defn}
A \textbf{homotopy of paths} in $X$ is a family
$f_t:I\to X$, with $0\leq t\leq0$, such that the
associated function
$$
\begin{aligned}
F:I\times I&\to X\\
(s,t)&\mapsto f_t(s)
\end{aligned}
$$
is continuous.

We say that the homotopy $F:I\times I\to X$ is
\textbf{relative to the endpoints} if
$$
F(0,t)=x_0\qquad\forall t\in I
\qquad
F(0,t)=x_1\qquad\forall t\in I
$$
for $x_0,x_1\in X$.

Finally, two paths $f_0$ and $f_1$ in $X$ such that
$f_0(0)=f_1(0)$ and $f_0(1)=f_1(1)$ are
\textbf{homotopic rel} $\mathbf{0,1}$ if there is a
homotopy between them. This is denoted
$f_0\simeq f_1\rel0,1$.
\end{defn}

\begin{teo}
The relation
$f_1\sim{}f_2\iff f_1\simeq f_2\rel0,1$ is an
equivalence relation on the set of paths in $X$, which we
will denote by $\Omega(X):=\{f:I\to X\}$.
\end{teo}

\section{Definition of the fundamental group}

On our way to defining the fundamental group, we first
define an operation between paths.

We denote by $[f]$ the homotopy class rel $0,1$ of
$f\in\Omega(X)$.

\begin{defn}
Given $f,g:I\to X$ such that $f(1)=g(0)$, we can define
the \textbf{concatenation of paths}
$$
\begin{aligned}
f&\cdot g:I\to X\\
f\cdot g(s)&=\begin{cases}
f(2s)\qquad s\in[0,1/2]\\
g(2s-1)\qquad s\in[1/2,1].
\end{cases}
\end{aligned}
$$
\end{defn}

\begin{obs}
If $f\simeq f_1\rel0,1$ and $g\simeq g_1\rel0,1$, then
$f\cdot g\simeq f_1\cdot g_1\rel0,1$.

In other words, if $[f]=[f_1]$ and $[g]=[g_1]$, then
$[f\cdot g]=[f_1\cdot g]$.
\end{obs}

\begin{defn}
We say that $f:I\to X$ is a \textbf{loop based at}
$x_0\in X$ if $f(0)=f(1)=x_0$. We denote by
$$
\Omega(X,x_0)=\{f:I\to X:f(0)=f(1)=x_0\}
$$
the set of such loops.
\end{defn}

The equivalence relation $\rel0,1$ restricts correctly to
$\Omega(X,x_0)\subseteq\Omega(X)$, so we can finally define:

\begin{defn}
The \textbf{fundamental group} of $X$ based at $x_0$ is the
set
$$
\pi_1(X,x_0)=\Omega(X,x_0)/\rel0,1.
$$
\end{defn}

\begin{teo}
The fundamental group is a group with the concatenation
operation $[f][g]=[f\cdot g]$.
\end{teo}

\begin{proof}
We have to check that:
\begin{itemize}
\item the product is well defined;
\item there is an identity, the constant path
$c_{x_0}:I\to X$ given by $c_{x_0}(t)=x_0$ for every
$t\in I$;
\item there are inverses;
\item there is associativity.
\end{itemize}
\end{proof}

\section{Change of base point}

What happens when we change the base point? Does the
fundamental group change? It is easy to note that if the
points lie in different path-connected components, the
fundamental groups may very well change. But what happens
if they lie in the same path-connected component?

\begin{teo}\label{teo:cambioptobase}
If $x_0,x_1\in X$ can be connected by a path, then
$\pi_1(X,x_0)\cong\pi_1(X,x_1)$.
\end{teo}

\section{The fundamental group of the circle}

Let us think of $S^1=\{z\in\C:|z|=1\}$.

\begin{teo}
$\pi_1(S^1,x_0)\cong\Z$ for every $x_0\in S^1$.
\end{teo}

\section{Consequences of the fundamental group of the circle}

\section{The fundamental theorem of algebra}

\section{Brouwer's fixed point theorem in dimension 2}

\begin{teo}
Every function $f:D^2\to D^2$ has a fixed point.
\end{teo}

\section{The Borsuk--Ulam theorem in dimension 2}

\begin{teo}
For every continuous function $f:S^2\to\R^2$ there exists
$x\in S^2$ such that $f(x)=f(-x)$.
\end{teo}

\section{The fundamental group of a product}

\begin{teo}
Let $X$ and $Y$ be topological spaces and let $x_0\in X$,
$y_0\in Y$. Then
$$
\pi_1(X\times Y,(x_0,y_0))
\cong
\pi_1(X,x_0)\times\pi_1(Y,y_0).
$$
\end{teo}

\section{Induced homomorphisms (functoriality)}

First let us establish some notation. If $f:X\to Y$,
$x_0\in X$, $y_0\in Y$, and $f(x_0)=y_0$, we write
$f:(X,x_0)\to(Y,y_0)$.

\begin{teo}[The fundamental group is a functor
$\pi_1:\Top^*\to\Grp$]
The function $f:(X,x_0)\to(Y,y_0)$ induces a homomorphism
$$
f_*:\pi_1(X,x_0)\to\pi_1(Y,y_0)
$$
such that
\begin{enumerate}
\item if $g:(Y,y_0)\to(Z,z_0)$, then
$(g\circ f)_*=g_*\circ f_*$;
\item $(Id_X)_*=Id_{\pi_1(X,x_0)}$.
\end{enumerate}
\end{teo}

\section{The fundamental group of the $n$-sphere}

\begin{prop}
$\pi_1(S^n)=0$ if $n\geq2$.
\end{prop}

This is proved using the following lemma.

\begin{lema}\label{lema:homotopia}
Let $X=\bigcup_{\alpha\in I} A_\alpha$ be such that
\begin{itemize}
\item the $A_\alpha$ are open sets, i.e. they are an open
cover;
\item there exists $x_0\in\bigcap A_\alpha$;
\item $A_\alpha$ is path connected;
\item $A_\alpha\cap A_\beta$ is also path connected for all
$\alpha,\beta$.
\end{itemize}
Then every loop $\gamma$ in $X$ is homotopic rel $0,1$ to a
product, i.e. concatenation, of loops of the form
$\gamma\cong\gamma_1\cdots\gamma_n$ such that each
$\gamma_i$ is completely contained in some $A_\alpha$.
\end{lema}

\section{$\R^2$ is not homeomorphic to $\R^n$ if $n\neq2$}

\begin{teo}
$\R^2$ is not homeomorphic to $\R^n$ if $n\neq2$.
\end{teo}

\begin{obs}
If $x\in\R^n$, then $\R^n-\{x\}$ is homeomorphic to
$S^{n-1}\times\R$.
\end{obs}

\begin{teo}
$\R^n\cong\R^m\iff n=m$.
\end{teo}

\section{Homotopy}
\label{def:func-homot}

Two functions $f,g:X\to Y$ are \textbf{homotopic} if there
exists a function, which we call a \textbf{homotopy}, of the
form $H:X\times I\to Y$ such that
$$
H(x,0)=f(x)\qquad\forall x\in X,
\qquad
H(x,1)=g(x)\qquad\forall x\in X,
$$
and this is denoted by $f\simeq g$.

If we have $A\subseteq X$ and $B\subseteq Y$, and $f$ is
such that $f(A)\subseteq B$, we write
$$
f:(X,A)\to(Y,B).
$$
A homotopy between functions of this kind is a function
$H:X\times I\to Y$ such that all the functions obtained by
changing the time parameter still send $A$ to $B$. That is,
we have functions
$$
H(\_,t):(X,A)\to(Y,B)\qquad\forall t\in I
$$
such that
$$
H(x,0)=f(x)\qquad\forall x\in X,
\qquad
H(x,1)=g(x)\qquad\forall x\in X.
$$

In particular, if $A=\{x_0\}$ and $B=\{y_0\}$, just as
before, we write
$$
f:(X,x_0)\to(Y,y_0).
$$
In this case, we say that a \textbf{base-point-preserving
homotopy} is a function $H:X\times I\to Y$ such that
$$
H(x_0,t)=y_0\qquad\forall t\in I,
$$
$$
H(x,0)=f(x)\qquad\forall x\in X,
\qquad
H(x,1)=g(x)\qquad\forall x\in X.
$$
We will simply say that $f$ and $g$ are homotopic,
specifying when necessary that the homotopy preserves the
base point.

\begin{prop}
If $f,g:(X,x_0)\to(Y,y_0)$ are homotopic, preserving the
base point, then they induce the same homomorphism on
fundamental groups, that is,
$$
f_*=g_*:\pi_1(X,x_0)\to\pi_1(Y,y_0).
$$
\end{prop}

\begin{defn}
Two spaces $X$ and $Y$ are \textbf{homotopy equivalent} if
there exist two functions, called \textbf{homotopy
equivalences}, $f:X\to Y$ and $g:Y\to X$ such that
$g\circ f\simeq Id_X$ and $f\circ g\simeq Id_Y$. This is
denoted by $X\simeq Y$.
\end{defn}

\begin{coro}\label{1.2.1}
If $f:X\to Y$ is a homotopy equivalence, then the induced
function $f_*:\pi_1(X,x_0)\to\pi_1(Y,y_0)$ is an
isomorphism.
\end{coro}

\begin{proof}
Suppose that $g:Y\to X$ is as in the preceding definition.
It is enough to show that $f_*$ and $g_*$ are inverses of one
another. And they are, because
$Id_{\pi_1(X,x_0)}=(g f)_*=g_*f_*$ and similarly
$Id_{\pi_1(Y,y_0)}=f_*g_*$.
\end{proof}

\begin{ejer}[Hatcher 0.11]\label{hatcher:0.11}
\textit{One can use different functions to go and come
back.} A function $f:X\to Y$ is a homotopy equivalence if
there exist $g,h:Y\to X$ such that $fg\simeq Id_Y$ and
$gf\simeq Id_X$. More generally, $f$ is a homotopy
equivalence if $fg$ and $hf$ are homotopy equivalences.
\end{ejer}

We add one last property, for which we define two notions
that are used throughout the text.

\begin{defn}\label{def:retracto}
A \textbf{retraction} of a space $X$ onto a subspace $A$ is a
function $f:X\to A$ such that $f|_A=id_A$. A
\textbf{deformation retract} of $X$ onto $A$ is a homotopy
$f_t:X\to X$ such that $f_0=id_A$ and $f_1(X)=A$, and also
$f_t|_A=id_A$ for every $t$.
\end{defn}

\begin{prop}
If a space $X$ retracts onto a space $A$, then the
homomorphism $i_*:\pi_1(A,x_0)\to\pi_1(X,x_0)$ induced by
the inclusion $A\hookrightarrow X$ is injective. If $A$ is a
deformation retract of $X$, then $i_*$ is an isomorphism.
\end{prop}

\chapter{The Van Kampen theorem}

\section{Free groups}

\section{Free products}

In group theory, we normally define the direct product by
thinking of something like
$$
\xymatrix{
& G\times H \\
G_\alpha \ar[ur] && H \ar[ul]
}
$$
where the product is commutative, that is,
$$
(g,1)(1,h)=(1,h)(g,1).
$$
Now we define a product called the \textbf{free product}, in
which the elements will not commute.

\begin{defn}
Let $\{G_\alpha\}_{\alpha\in I}$ be a collection of groups.
As a set, the free product $\ast_{\alpha\in I}G_\alpha$
consists of the words $g_1g_2\ldots g_m$, for
$m\in\N\cup\{0\}$, where each $g_i$ lies in some
$G_\alpha$. Also, $g_i\neq1$, and $g_i,g_{i+1}$ always
belong to different $G_\alpha$, that is, $g_i\in G_\alpha$
and $g_{i+1}\in G_\beta$ with $\alpha\neq\beta$. In this
case, we say that $g_1\ldots g_m$ is a
\textbf{reduced word}.

The binary operation on $\ast_{\alpha\in I}$ is given by
concatenation.
\end{defn}

\begin{obs}
Each of the $G_\alpha$ is contained in
$\ast_{\alpha\in I}G_\alpha$, since its elements appear as
one-letter words. In symbols:
$$
\begin{aligned}
G&\hookrightarrow\ast_{\alpha\in I}G_\alpha\\
g&\mapsto g.
\end{aligned}
$$
\end{obs}

\begin{teo}[Universal property of the free product]
Given homomorphisms $\varphi_\alpha:G_\alpha\to H$, there
exists a unique homomorphism
$\varphi:\ast_{\alpha\in I}G_\alpha\to H$ extending the
$\varphi_\alpha$.
\end{teo}

\begin{ejem}
If $F_n$ and $F_m$ are free groups generated by $n$ and $m$
elements, then $F_n*F_m\approx F_{n+m}$.
\end{ejem}

\section{Normally generated subgroups}

\section{The Van Kampen theorem}

Imagine that we want to compute the fundamental group of a
space but do not know how. This happens very often. But it
turns out that we can view this space as the union of many
subspaces whose fundamental groups we know. This is the
setting of the Van Kampen theorem.

Take a topological space $X$, a base point $x_0\in X$, and
suppose that $X=\bigcup_{\alpha\in I}A_\alpha$ for certain
spaces $A_\alpha$ such that $x_0\in A_\alpha$ for every
$\alpha\in I$. Also suppose that the $A_\alpha$ are path
connected.

Since the $A_\alpha$ are subspaces, we can consider the
inclusion $A_\alpha\hookrightarrow X$ and the induced
homomorphisms on fundamental groups, denoted
$j_\alpha:\pi_1(A_\alpha,x_0)\to\pi_1(X,x_0)$. Then, by the
universal property of the free product, we have a
homomorphism from the free product of these groups to the
fundamental group of $X$:
$$
\ast_{\alpha\in I}\pi_1(A_\alpha,x_0)
\to
\pi_1(X,x_0).
$$
The theorem will tell us that this function is surjective.
Intuitively, it says that every loop in $X$ is homotopic to
another loop which is the concatenation of certain loops,
each completely contained in one of the $A_\alpha$.

Now think of a loop that lies in the intersection of two of
the $A_\alpha$. We can think of the loop as lying in, say,
$A_\alpha$, and its inverse as lying in $A_\beta$. Then, as
an element of the free product
$\ast_{\alpha\in I}\pi_1(A_\alpha,x_0)$, this loop is not the
identity. But after pushing it to $\pi_1(X,x_0)$, it is the
identity. All this is to say that the map we have built is
not injective.

Now suppose that for $\alpha,\beta\in I$ we have the
inclusion $A_\alpha\cap A_\beta\hookrightarrow A_\alpha$,
which induces a map
$$
i_{\alpha\beta}:\pi_1(A_\alpha\cap A_\beta)
\to
\pi_1(A_\alpha).
$$
The order of the subscripts is important, since the
codomain is given by the first symbol in the subscript.

Putting the preceding paragraphs together, the nontrivial
loop
$$
1\neq i_{\alpha\beta}[f]i_{\beta\alpha}[\bar f]
\in \ast_{\alpha\in I}A_\alpha
$$
can be pushed to $\pi_1(X,x_0)$. The Van Kampen theorem
will tell us that the kernel of the homomorphism, which we
already said is certainly not injective in general, is
generated precisely by elements of this kind.

Now we can state it.

\begin{teo}[Van Kampen]
If $X=\bigcup_{\alpha\in I}A_\alpha$, with each $A_\alpha$
path connected, $x_0\in A_\alpha$, and also
$A_\alpha\cap A_\beta$ path connected for all
$\alpha,\beta\in I$, then the homomorphism
$$
\Phi:\ast_{\alpha\in I}\pi_1(A_\alpha,x_0)
\to\pi_1(X,x_0)
$$
is surjective.

If moreover $A_\alpha\cap A_\beta\cap A_\gamma$ is path
connected for all $\alpha,\beta,\gamma\in I$, then
$$
\ker\Phi=
\left\langle\left\langle
 i_{\alpha\beta}(\omega)i_{\beta\alpha}(\omega)^{-1}
 :\omega\in\pi_1(A_\alpha\cap A_\beta)
\right\rangle\right\rangle.
$$
In particular,
$$
\pi_1(X,x_0)\cong
\ast_{\alpha\in I}\pi_1(A_\alpha,x_0)/\ker\Phi.
$$
\end{teo}

\section{The fundamental group of a wedge of spaces}
\label{sec:grp-fund-cuña}

\textcolor{red}{I did not find
\href{https://www.youtube.com/watch?v=9P6n__Njlz8}{this video}
on the course page.}

Consider a collection of topological spaces $X_\alpha$ with
$\alpha\in I$ and choose a point in each space, say
$x_\alpha\in X_\alpha$. We define the \textbf{wedge} of these
spaces by
$$
\bigvee_{\alpha\in I}X_\alpha
:=
\bigsqcup_{\alpha\in I}X_\alpha
\Big/\{x_\alpha\sim{}x_\beta,\forall\alpha,\beta\in I\}.
$$
Assume also that each $x_\alpha$ is a deformation retract
of some neighbourhood $U_\alpha\subseteq X_\alpha$.

Then we have:

\begin{teo}
Under these hypotheses, the homomorphism appearing in the
Van Kampen theorem,
$$
\Phi:\ast_{\alpha\in I}\pi_1(X_\alpha,x_0)
\to
\pi_1\left(\bigvee_{\alpha\in I}X_\alpha\right),
$$
is an isomorphism.
\end{teo}

\section{Every group is the fundamental group of some space}

\section{Presentations of groups}

Consider a group $G$ and a generating set $S\subseteq G$,
that is, such that $\langle S\rangle=G$. This is equivalent
to saying that every element of $G$ can be seen as a product
of elements of $S$.

Also consider the free group generated by $S$, say $F_S$,
and use the universal property of free groups to see that
there exists a unique surjective homomorphism from $F_S$ to
$G$ extending the inclusion of $S$ in $G$.

We have proved:

\begin{teo}
Given a group $G$, there exists a free group $F_S$ which
surjects onto $G$. In particular, by the first isomorphism
theorem, every group is a quotient of a free group.
\end{teo}

If in addition we have a set $R\subseteq\ker\varphi$ which
normally generates the kernel, that is,
$\langle\langle R\rangle\rangle=\ker\varphi$, then
$$
G\cong F_S/\ker\varphi
\cong F_S/\langle\langle R\rangle\rangle,
$$
and we say that
$$
G=\langle S,R\rangle
$$
is a \textbf{presentation} of $G$.

\begin{ejems}
\begin{enumerate}
\item S
\end{enumerate}
\end{ejems}

\begin{prop}
If $G=\langle S,R\rangle$ is a group presentation, $K$ is
another group, and $f:S\to K$ is a function such that for
any $s_1,\ldots,s_m\in R$, with $s_i\in S$,
$$
f(s_1)\ldots f(s_m)=1,
$$
then there exists a unique homomorphism $\bar f:G\to K$
such that the corresponding diagram commutes.
\end{prop}

That is, to define a homomorphism out of a group given by a
presentation, we only have to make sure that the relations
are respected.

\section{Every group is isomorphic to the fundamental group
of some space}

\begin{lema}
Let $X$ be a topological space, let $x_0\in X$ be a base
point, and let $f:S^1\to X$ be such that $f(1)=x_0$. For the
adjunction space $Y=X\cup_f D^2$,
$$
\pi_1(Y,x_0)\cong
\pi_1(X,x_0)\big/\langle\langle [f]\rangle\rangle,
$$
where we view $f$ as a loop in the fundamental group.
\end{lema}

This is called ``killing'' a loop.

\begin{lema}
Let $X$ be a topological space, let $x_0\in X$ be a base
point, and let $f_\alpha:S^1\to X$, with $\alpha\in I$, be
such that $f(1)=x_0$. For the adjunction space
$Y=X\cup_{\{f_\alpha\}}\{D^2\}$,
$$
\pi_1(Y,x_0)\cong
\pi_1(X,x_0)
\big/
\langle\langle [f_\alpha]:\alpha\in I\rangle\rangle,
$$
where we view $f$ as a loop in the fundamental group.
\end{lema}

And now:

\begin{teo}
For every group $G$ there exists a topological space $X$
such that $\pi_1(X,x_0)\cong G$.
\end{teo}

\begin{proof}
First take a presentation
$G=\langle S,R\rangle=F_S/\langle\langle R\rangle\rangle$,
and then take the wedge of as many circles as there are
elements in $S$, say $X=\bigvee_{\alpha\in S}$. We know that
$\pi_1(X,x_0)\approx *_{\alpha\in S}\Z\approx F_S$.

Now we have to kill as many loops as there are relations. A
relation is a word in the generators, say
$r=s_1^{\varepsilon_1}s_2^{\varepsilon_2}\ldots
s_m^{\varepsilon_m}$ with $s_i\in S$ and
$\varepsilon_i=\pm1$. This word essentially determines a
path in the wedge of circles, so we can glue in one disk for
each relation, and by the preceding lemma we are done.
\end{proof}

\section{The fundamental group does not detect cells of
higher dimension than 2}

To formalize the title of this section, we need:
\begin{itemize}
\item a connected topological space $X$;
\item a cell $e^n$ of dimension $n>2$, that is, a space
homeomorphic to the closed $n$-disk;
\item an attaching map $f:S^{n-1}\to X$;
\item the adjunction space $Y\cup_f e^n$;
\item a point $x_0\in\img(f)$.
\end{itemize}
Then we have:

\begin{teo}
The inclusion $i:X\hookleftarrow Y$ induces an isomorphism
$i_*:\pi_1(X,x_0)\to\pi_1(Y,y_0)$.
\end{teo}

That is, when we glue a cell of dimension greater than 2 to
a space, the fundamental group stays the same.

\chapter{Exercises}

\begin{ejer}
For a path-connected space $X$, show that $\pi_1(X)$ is
abelian if and only if the change-of-base-point
homomorphisms $\beta_h$ depend only on the endpoint of $h$.
\end{ejer}

\begin{proof}[Solution]
For the forward direction, take
$[\gamma],[\delta]\in\pi_1(X,x_0)$. We know that for any
paths $h,g$ ending at $x_0$, the homomorphisms $\beta_h$ and
$\beta_g$ are equal, so
$$
[h\cdot \gamma\cdot\delta\cdot\bar h]
=
[g\cdot\gamma\cdot\delta\cdot\bar g].
$$
Take $h=\delta$ and $g=c_{x_0}$. Then
$\delta\cdot\gamma\cdot\delta\cdot\bar\delta
\simeq\gamma\cdot\delta$.

For the converse, take two paths $h$ and $g$ ending at
$x_0$, and take $[\gamma]\in\pi_1(X,x_0)$. Now note that
$$
[h\cdot\gamma\cdot\bar h]
=[h\cdot\gamma\cdot \bar g\cdot g\cdot\bar h]
=[h\cdot\gamma\cdot \bar g][g\cdot\bar h]
=[\bar g\cdot\bar h][h\cdot\gamma\cdot \bar g]
=[g\cdot\gamma\cdot \bar g].
$$
\end{proof}

\begin{ejer}
Suppose that $f_t:X\to X$ is a homotopy such that $f_0$ and
$f_1$ are the identity. Show that for every $x_0\in X$, the
loop $f_t(x_0)$ represents an element in the center of
$\pi_1(X,x_0)$.
\end{ejer}

\begin{proof}[Solution]
Let $[f]$ be the equivalence class of the loop $f_t(x_0)$,
and let $[\gamma]$ be any element of $\pi_1(X,x_0)$. We want
to prove that $f\cdot\gamma\simeq\gamma\cdot f$.

One can think of loops like $f_t(x_0)$ but starting at any
point of the space $X$. In particular we have the family of
loops $f_t(\gamma(s))$. The omitted diagram represents the
domain of this function. Of course
$f_0(\gamma(s))=\gamma(s)=f_1(\gamma(s))$, and
$f_t(\gamma(0))=f_t(x_0)=f_t(\gamma(1))$. Thus the left and
top edges represent $f\cdot\gamma$, and the bottom and right
edges represent $\gamma\cdot f$. Any path inside the square
joining the lower-left corner to the upper-right corner
represents a path homotopic to either of the above.
\end{proof}

\begin{ejer}
Use Lemma \ref{lema:homotopia} to show that if $X$ is
obtained from a path-connected space by adjoining a cell
$e^n$ with $n\geq2$, then the inclusion $A\hookrightarrow X$
induces a surjective function on $\pi_1$.
\end{ejer}

\begin{proof}[Solution]
Lemma \ref{lema:homotopia} ensures that if we find two open
sets $A_1$ and $A_2$ containing the base point such that
their union is the whole space and their intersection is path
connected, then every loop in the space can be seen as the
product of two loops, each contained in some $A_i$.

Take $A_1=e^n$ and $A_2=X-c$, where $c$ is a point in the
interior of $e^n$. If we choose a base point $x_0$ also in
the interior of $e^n$, we have the conditions for the lemma.
Then every loop in $X$ can be seen as a loop only in $A_2$,
since $e^n$ is contractible. Then we must use the
change-of-base-point theorem to connect $x_0$ with a
$\tilde{x}_0$ lying in $X\backslash e^n$. Finally, since $A$
is a deformation retract of $A_2$, their fundamental groups
are isomorphic, so we can think of our loop as lying in $A$.
\end{proof}

\part{Covering spaces}

\chapter{Covering spaces}

\section{Definition of covering space}

In the proof of the fundamental group of the circle we had
already used the idea of a covering space. Here is the
general definition.

\begin{defn}
A \textbf{covering space} of $X$ is a space $\tilde X$
together with a function
$$
p:\tilde X\to X
$$
such that each point $x\in X$ has an open neighbourhood
$U\subseteq X$ such that $p^{-1}(U)$ is a disjoint union of
open subsets of $\tilde X$, and each of these projects
homeomorphically onto $U$ via $p$.
\end{defn}

% Covering-space illustration omitted.

We introduce other terms:
\begin{itemize}
\item the function $p$ is called the \textbf{covering
projection};
\item $U$ is called an \textbf{evenly covered neighbourhood};
\item the open subsets of $p^{-1}(U)$ which project
homeomorphically to $U$ are called the \textbf{sheets} of
$\tilde X$ over $U$;
\item the set $p^{-1}(x)$ is called the \textbf{fiber} of
$x$ under $p$.
\end{itemize}

\begin{obs}\leavevmode
\begin{itemize}
\item If $U$ is connected, then the sheets over $U$ are the
connected components of $p^{-1}(U)$.
\item The number of sheets over $U$ is equal to the
cardinality of $p^{-1}(x)$.
\item This number is \textit{locally constant}, that is, it
is the same for any point.
\end{itemize}
\end{obs}

\begin{ejems}\leavevmode
\begin{enumerate}
\item The identity $Id:X\to X$.
\item The natural projection from the disjoint union of
copies of a space, that is,
$p:\bigsqcup_{\alpha\in I}X\to X$.
\item For any two coverings $p:\tilde X\to X$ and
$q:\hat X\to X$, the projection from the disjoint union
$r:\tilde X\sqcup\hat X\to X$, defined in the natural way,
is also a covering. This is to say that we are interested in
coverings of connected spaces whose covering spaces are
connected.
\item $p:\R\to S^1$, as in the proof of the fundamental
group of the circle.
\item For $n\in\N$, the function $p:S^1\to S^1$ given by
$p(z)=z^n$.
\end{enumerate}
\end{ejems}

\section{Homotopy lifting for coverings}

\begin{defn}
A \textbf{lift} of a function $f:Y\to X$ with respect to a
covering $p:\tilde X\to X$ is a function
$\tilde f:Y\to\tilde X$ such that $p\tilde f=f$.
Equivalently, the corresponding triangle commutes.
\end{defn}

\begin{teo}[Homotopy lifting property for coverings]
\label{teo-lhc}
Given a covering $p:\tilde X\to X$, a homotopy
$h_t:Y\to X$, and a lift $\tilde f_0:Y\to\tilde X$ of
$f_0$, there exists a unique homotopy
$\tilde f_t:Y\to\tilde X$ lifting each of the $f_t$.
\end{teo}

\section{The induced homomorphism of the covering projection}

Given a covering projection $p:\tilde X\to X$, what
properties does the induced map
$$
p_*:\pi_1(\tilde X,\tilde x_0)\to (X,x_0)
$$
have, where $p(\tilde x_0)=x_0$? It will turn out to be
injective.

To prove this, we use the homotopy lifting property for
coverings. First we look at a couple of simple cases of how
this theorem is used.

\begin{prop}[Path lifting property]
A path has a unique lift once the image of its initial point
has been defined. To see this, choose the case in which $Y$
consists of a single point, so that a homotopy of
$Y\times I$ is a path in $X$.
\end{prop}

\begin{obs}
The lift of a constant path is a constant path.
\end{obs}

\begin{obs}
If $Y=I$ and $H:I\times I$ is a homotopy rel $0,1$, then the
lift $\tilde H$ is also a homotopy rel $0,1$.
\end{obs}

Now:

\begin{prop}
Given a covering projection $p:\tilde X\to X$, the induced
map $p_*:\pi_1(\tilde X,\tilde x_0)\to (X,x_0)$ is
injective. Moreover,
$p_*(\pi_1(\tilde X,\tilde x_0))\leq\pi_1(X,x_0)$ consists
of the loops based at $x_0$ which lift to loops in
$\tilde X$ based at $\tilde x_0$.
\end{prop}

\section{The number of sheets of a covering and the
fundamental group}

\begin{prop}
Let $p:\tilde X\to X$ be a covering with $\tilde X$ and $X$
path connected. Then the number of sheets of the covering is
equal to the index of $p_*(\pi_1(\tilde X,\tilde x_0))$ in
$\pi_1(X,x_0)$.
\end{prop}

\section{The lifting criterion for functions}

Consider a covering $p:\tilde X\to X$ and an arbitrary
function $f:Y\to X$. Under what conditions is it possible to
lift this function to the covering? We want the usual
triangle to commute. Such functions do not always exist. To
know when they do, first let us see:

\begin{defn}
We say that a topological space $Y$ is \textbf{locally path
connected} if for every point $y\in Y$ and every
neighbourhood $U$ of $y$, there exists
$y\in V\subseteq U$ with $V$ path connected.
\end{defn}

\begin{prop}
Let $p:(\tilde X,\tilde x_0)\to (X,x_0)$ be a covering and
$f:(Y,y_0)\to (X,x_0)$ a function, with $Y$ path connected
and locally path connected. Then there exists a lift
$\tilde f:(Y,y_0)\to(\tilde X,x_0)$ if and only if
$$
f_*\pi_1(Y,y_0)\subseteq
p_*\pi_1(\tilde X,\tilde x_0).
$$
\end{prop}

\section{Uniqueness of lifts of functions}

Under what conditions is a lift unique? When two lifts agree
at one point, they have to be equal, provided the domain is
path connected.

\begin{prop}
Let $p:\tilde X\to X$ be a covering and $f:Y\to X$ a
function. If two lifts $\tilde f_1,\tilde f_2:Y\to\tilde X$
of $f$ agree at one point and $Y$ is connected, then
$\tilde f_1=\tilde f_2$.
\end{prop}

\section{The universal covering}

In this section we try to classify all coverings a space can
have. The first question we ask is in what cases a covering
has a covering whose fundamental group is trivial. Such a
covering will be called the \textbf{universal covering}; it
is universal in the sense that it will cover any other
covering.

Suppose we have a space $X$ which
\begin{itemize}
\item is path connected;
\item is locally path connected.
\end{itemize}
Now take a covering $p:\tilde X\to X$ as above, simply
connected. When we take a point $p\in X$ and an evenly
covered neighbourhood, any loop contained in that
neighbourhood can be lifted to one of its preimages. Since
the covering $\tilde X$ is simply connected, the lifted loop
is contractible, and this homotopy projects to a homotopy
from the original loop to the constant loop.

\begin{defn}
A space is \textbf{semilocally simply connected} if for every
$x\in X$ there exists a neighbourhood $U$ such that every
loop contained in $U$ retracts to the constant loop $c_x$ in
$X$. Equivalently, the homomorphism induced by inclusion
$\pi_1(U,x_0)\to(X,x_0)$ is constant.
\end{defn}

Thus we have already shown that if $X$ has a universal
covering, then it is semilocally simply connected. Now we
want to see that the converse is also true, when $X$ is also
path connected and locally path connected.

\begin{prop}
A space $Y$ is simply connected if and only if for every
$x,y\in Y$ there exists a unique homotopy class rel $0,1$ of
paths joining $x$ to $y$.

This implies that there is a bijection between the points of
$Y$ and the equivalence classes of paths starting at $x_0$,
that is,
$$
Y\leftrightarrow \{[\gamma]:\gamma(0)=x_0\}.
$$

Now consider a covering $p:\tilde X\to X$ with $\tilde X$
simply connected and a point $\tilde x\in\tilde X$. Denote
$p(\tilde x_0)=x_0\in X$. Then, by what was just said, the
points in the covering space $\tilde X$ are in bijection with
the paths starting at $\tilde x$, and projecting by $p$, we
arrive at the paths starting at $x_0$. Using the homotopy
lifting property, one verifies that this correspondence is
bijective:
$$
\{[\gamma]:\gamma(0)=\tilde x_0\}
\leftrightarrow
\{[\beta]:\beta(0)=x_0\}.
$$
\end{prop}

Now:

\begin{teo}
Suppose that $X$ is a path-connected, locally path-connected,
and semilocally simply connected topological space, and take
a base point $x_0\in X$. The space
$$
\tilde X:=\{[\gamma]:\gamma(0)=x_0\},
$$
together with the function
$$
\begin{aligned}
p:\tilde X&\to X\\
[\gamma]&\mapsto\gamma(1),
\end{aligned}
$$
is a universal covering of $X$. That is:
\begin{itemize}
\item $\tilde X$ has a topology;
\item $p$ is continuous;
\item $p$ is a covering;
\item $\tilde X$ is simply connected.
\end{itemize}
\end{teo}

\section{The classification theorem for coverings}

In this section we will see, roughly speaking, that there is
a correspondence between universal coverings, up to
isomorphism, and the subgroups of the fundamental group.

\begin{prop}
Let $X$ be path connected, locally path connected, and
semilocally simply connected. Then for every
$H\leq\pi_1(X,x_0)$ there exists a covering $p:X_H\to X$
such that
$$
p_*\left(\pi_1(X_H,\tilde x_0)\right)=H.
$$
\end{prop}

\begin{proof}
Using the construction of the universal covering, take an
equivalence relation on paths which identifies paths whose
concatenation lies in $H$. The quotient gives the desired
covering.
\end{proof}

\begin{defn}
A covering $p:\tilde X\to X$ is \textbf{normal} if
$p_*(\pi_1(\tilde X,\tilde x_0))$ is a normal subgroup of
$\pi_1(X,x_0)$.
\end{defn}

\begin{prop}
If $p:\tilde X\to X$ and $q:\hat X\to X$ are coverings of
$X$, with all spaces path connected and locally path
connected, then a lift $\tilde X\to\hat X$ over $X$ exists if
and only if
$$
p_*\pi_1(\tilde X)\subseteq q_*\pi_1(\hat X).
$$
\end{prop}

\begin{teo}[Classification of coverings]
Let $X$ be path connected, locally path connected, and
semilocally simply connected. There is a correspondence
between isomorphism classes of path-connected coverings of
$X$ and conjugacy classes of subgroups of $\pi_1(X,x_0)$.
\end{teo}

\begin{teo}
The normal coverings correspond to the normal subgroups of
$\pi_1(X,x_0)$.
\end{teo}

\section{Deck transformations}

A \textbf{deck transformation} of a covering $p:\tilde X\to X$
is a homeomorphism $h:\tilde X\to\tilde X$ such that
$p\circ h=p$. In other words, it is an automorphism of the
covering over $X$.

The deck transformations form a group under composition.

\section{Group actions and coverings}\label{sec:acciones}

\begin{defn}
An \textbf{action} of a group $G$ on a space $X$ is a
homomorphism from $G$ to the group of homeomorphisms of $X$.
We write $G\curvearrowright X$.

The \textbf{orbit} of a point $x\in X$ is the set
$\Orb\{g(x):g\in G\}$.
\end{defn}

Notice that the relation $x\sim{}y\iff y\in\Orb(x)$ is an
equivalence relation, so we can think of the quotient space
$X/\sim{}:=X/G$. We ask: given a group action
$G\curvearrowright Y$, when is $Y/G$ a covering?

\begin{defn}
We say that $G\curvearrowright Y$ is a \textbf{covering
action} if either of the following equivalent conditions
holds:
\begin{itemize}
\item every point $y\in Y$ has a neighbourhood $U$ such that
$(g_1U\cap g_2U)\neq\emptyset\implies g_1=g_2$;
\item every point $y\in Y$ has a neighbourhood $U$ such that
$gU\cap U\neq\emptyset\implies g=1_G$.
\end{itemize}
In particular, these conditions imply that the action is
\textbf{free}, meaning that if $g(y)=y$ for any $g\in G$ and
$y\in Y$, then $g=1_G$.
\end{defn}

\begin{prop}
Given a covering action $G\curvearrowright Y$,
\begin{enumerate}
\item the quotient projection $p:Y\to Y/G$ given by
$p(y)=Gy$ is a normal covering;
\item if $Y$ is path connected, then $G$ is the group of deck
transformations of $p$;
\item if $Y$ is path connected and locally path connected,
then
$$
G\cong\pi_1(Y/G,y_0)/p_*\pi_1(Y,\tilde y_0).
$$
\end{enumerate}
\end{prop}

\begin{coro}
If moreover $Y$ is simply connected, then by (b) and (c),
$G=G(\tilde Y)\cong\pi_1(Y/G,y_0)$.
\end{coro}

\chapter{Exercises}

\begin{ejer}
\textbf{Let $S$ be the 2-sphere and $T$ the 2-torus. Compute
the fundamental group $G$ of $X:=S\vee T$, construct the
universal covering $\widetilde X\to X$, and describe the
action by deck transformations of $G$ on $\widetilde X$.}
\begin{proof}[Solution]
By the Van Kampen theorem,
$\pi_1(S\vee T,x_0)=\Z\times\Z$. The universal covering is
$\R^2$ with copies of $S^2$ glued at each point with integer
coordinates, and the action by deck transformations is just
the collection of translations by vectors with integer
coordinates.
\end{proof}
\end{ejer}

\begin{ejer}
\textbf{Let $f:X\to Y$ be a covering. Show that if $Y$ is a
CW complex, then so is $X$.}
\end{ejer}

\begin{proof}[Solution]
First define $X^0:=f^{-1}(Y^0)$. The idea is to use the
lifting theorem for functions on the attaching maps.
However, this cannot be done to construct the 1-skeleton of
$X$, because $S^0$ is not connected. What we can do is take
the characteristic map, which lifts because $D^1_\alpha$ is
contractible. With this we can define the attaching map
$\tilde\varphi_\alpha:S^0\to X^0$ directly by
$\tilde\varphi_\alpha(-1)=\tilde\Phi_\alpha(0)$ and
$\tilde\varphi_\alpha(1)=\tilde\Phi_\alpha(1)$, which works
because the characteristic map extends the attaching map.

The same method works to construct the 2-skeleton, and the
other attaching maps can be lifted directly because $S^n$ is
simply connected for $n>1$.

If $Y$ is finite-dimensional, we know that $Y=Y^n$.
\textit{A little bit is missing.}
\end{proof}

\begin{ejer}
\textbf{Show that for every $g\geq2$ there exists a covering
$S_g\to S_2$, that is, every surface of genus 2 is covered
by any surface of genus $g\geq2$. Also show that the torus is
not covered by any closed surface of genus $g\geq2$.}
\end{ejer}

\begin{proof}
Here we should visualize $S_g$ as a doughnut-flower with a
hole in the center and $g-1$ petals. In this case, the cyclic
group of order $g-1$ is a covering action on $S_g$, rotating
around the central hole. The resulting quotient space is just
$S_2$, since all the petals are identified.

Suppose that there exists a covering $p:S_g\to T$ with
$g\geq2$. We know that the induced homomorphism $p_*$ is
injective, so $p_*\pi_1(S_g)\approx\pi_1(S_g)$, which by Van
Kampen is of the form
$$
\langle a_1,b_1,\ldots,a_g,b_g
\mid
 a_1b_1a_1^{-1}b_1^{-1}\ldots
 a_gb_ga_g^{-1}b_g^{-1}\rangle,
$$
which is not commutative. But every subgroup of
$\pi_1(T)\approx\Z\times\Z$ is commutative.
\end{proof}

\begin{ejer}[Hatcher, 1.3.8]
\textbf{Let $\tilde X$ and $\tilde Y$ be simply connected
covering spaces of path-connected and simply connected
spaces $X$ and $Y$. Show that if $X\simeq Y$, then
$\tilde X\simeq\tilde Y$. Perhaps Exercise
\ref{hatcher:0.11} is useful.}
\end{ejer}

\begin{proof}[Solution]
Everything is in the corresponding commutative diagram. Here
$f$ and $g$ are homotopy equivalences between $X$ and $Y$.
The composition $pf$ lifts to a unique function $\tilde f$ by
the connectivity hypotheses, and similarly we find
$\tilde g$. Since everything commutes, $\tilde g\tilde f$ is
a lift of $gfp$.

Since $gf\simeq Id_X$, we have $gfp\simeq p$. This homotopy
lifts uniquely to a homotopy
$H:\tilde X\times I\to\tilde X$, since
$H(\underline{\hspace{5pt}},0)=\tilde g\tilde f$. Thus
$pH(\underline{\hspace{5pt}},1)=p$.

So $H(\underline{\hspace{5pt}},1):=h$ is a lift of $p$, just
like the identity. If we could see that they agree at one
point, then by uniqueness of lifting functions we would have
$h=Id_{\tilde X}$.

But we are not certain this occurs. What we do know is that,
since $\tilde X$ is simply connected, $p$ is a normal
covering. Thus the group of deck transformations acts
transitively on the fibers. Take a point $\xi\in\tilde X$. If
$h$ does not fix it, no matter; there exists a deck
transformation $u$ such that $u(h(\xi))=\xi$.

Now by uniqueness of lifts, $uh=Id_{\tilde X}$. Hence
$u\tilde g\tilde f\simeq Id_{\tilde X}$.

Exercise \ref{hatcher:0.11} ensures that if we find two
functions $\phi,\psi:\tilde X\to\tilde X$ such that
$\phi\tilde f\simeq Id_{\tilde X}$ and
$\tilde f\psi\simeq Id_{\tilde Y}$, then $\tilde f$ is a
homotopy equivalence and we are done. We have already found
$\phi$.

To find $\psi$, do the same reasoning, now requiring a deck
transformation $w$ such that
$\tilde f\tilde g w\simeq Id_{\tilde Y}$. That is, here one
must apply $w$ before the function obtained through the
homotopy, say $\bar h$. This is legal: instead of moving
$\bar h(\xi)$ back to $\xi$, beforehand we move $\xi$ by $w$
to a point in the preimage of $\xi$ under $\bar h$. Although
perhaps $\bar h$ is not surjective, it is enough to choose
$\xi\in\img\bar h$.
\end{proof}

\begin{ejer}
Show that if $X$ is path connected and locally path connected
and $\pi_1(X)$ is finite, then every function $f:X\to S^1$ is
homotopic to a constant.
\end{ejer}

\begin{proof}
Since $\pi_1(X)$ is finite and $f_*\pi_1(X)$ is a subgroup of
$\pi_1(S^1)\approx\Z$, and the only finite subgroup of $\Z$
is the trivial one, we have $f_*\pi_1(X)\approx0$. Thus there
exists a lift $\tilde f:X\to\R$ to the universal covering of
the circle.

Since $X$ is path connected, its image must be an interval,
which is contractible. Thus
$\img\tilde f\simeq\{\tilde x_0\}$ via a homotopy
$h:\img\tilde f\times I\to\img\tilde f$ such that
$h(x,0)=x$ and $h(x,1)=\tilde x_0$.

The function $g=p\circ h\circ\tilde f$ is a homotopy between
$f=p\circ\tilde f$ and the constant $x_0=p(\tilde x_0)$.
\end{proof}

\part{Homology}

\chapter{Homological algebra}

\section{Basic concepts}

In this chapter $R$ denotes an associative ring with unit,
not necessarily commutative. Usually we will think of one of
$\Z,\Q,\R$.

Recall that an $R$-module is basically a vector space, but
with scalars in $R$.

\begin{defn}
An \textbf{$R$-chain complex} is a sequence of $R$-modules
and homomorphisms
$$
\cdots \to C_p \xrightarrow{\partial_p} C_{p-1}
\xrightarrow{\partial_{p-1}} C_{p-2}\to\cdots
$$
such that $\partial_{p-1}\partial_p=0$ for every
$p\in\mathbb Z$, equivalently
$\text{img }\partial_p\subseteq\ker\partial_{p-1}$.
\end{defn}

\begin{defn}
A \textbf{morphism of $R$-chain complexes}
$(C_\bullet,\partial)\to(D_\bullet,\delta)$ is a sequence of
$R$-homomorphisms $C_p\xrightarrow[]{f_p}D_p$ such that the
usual square with the boundary maps commutes, that is,
$$
f_{p-1}\partial_p=\delta_pf_p
$$
for every $p\in\mathbb Z$.
\end{defn}

\begin{defn}
We say that $(D_\bullet,\delta)$ is a
\textbf{sub-chain-complex} of $(C_\bullet,\partial)$ if
$D_p\leq C_p$ for every $p\in\mathbb Z$ and
$\partial|_{D_p}=\delta_p$. The quotient
$(C_\bullet/D_\bullet,\partial)$ is the chain complex whose
terms are $C_p/D_p$, with boundary maps
$\partial_p/\delta_p([c])=[\partial_p(c)]$.
\end{defn}

\begin{defn}\leavevmode
\begin{itemize}
\item Elements of $C_p$ are called \textbf{$p$-dimensional
chains}.
\item Elements of $\ker\partial_p:=Z_p$ are called
\textbf{$p$-dimensional cycles}.
\item Elements of $\text{img }\partial_{p+1}:=F_p:=B_p$ are
called \textbf{$p$-dimensional boundaries}.
\end{itemize}
\end{defn}

\begin{defn}
The \textbf{$p$-th homology group} of $(C_\bullet,d)$ is
$$
H_p(C_{\dot{}}):=Z_p/B_p
=\ker\partial_p/\text{img }\partial_{p+1}.
$$
We say that two cycles $c$ and $c'$ are \textbf{homologous}
if $[c]=[c']\in H_p(C_\bullet)$.
\end{defn}

% Figure of homologous cycles omitted.

\begin{ejer}[Induced function]
If $(C_\bullet,\partial)\xrightarrow f(C'_\bullet,\partial')$
is a homomorphism, then $f(Z_p)\subseteq Z'_p$ and
$f(B_p)\subseteq B'_p$, so the induced function
$$
\begin{aligned}
\bar f_p:H_p(C_\bullet)&\to H_p(C_\bullet)\\
a+B_p&\mapsto f_p(a)+B'_p
\end{aligned}
$$
is well defined. If we also have a second homomorphism
$(C'_\bullet,\partial')\xrightarrow g(C''_\bullet,\partial'')$,
then $\overline{g\circ f}=\bar g\circ\bar f$. Finally,
$\overline{Id}_{C_p}=Id_{H_p(C)}$.
\end{ejer}

With this exercise we begin to see the functorial properties
of homology, although for now we will not go deeper into this
language.

\section{Exact sequences}

\begin{defn}
We say that the sequence
$$
\cdots \to C_p \xrightarrow{f_p} C_{p-1}
\xrightarrow{f_{p-1}} C_{p-2}\to\cdots
$$
is \textbf{exact at $C_p$} if
$\text{img }f_p=\ker f_{p-1}$. The sequence is
\textbf{exact} if it is exact at every $C_p$. This happens if
and only if $H_p(C_\bullet)=0$ for every $p\in\mathbb Z$.
\end{defn}

\begin{obs}\leavevmode
\begin{itemize}
\item The homology group measures how far the sequence is
from being exact.
\item The sequence may be ``finite'', that is, many modules
may be zero.
\end{itemize}
\end{obs}

\begin{defn}
An exact sequence of the form
$$
0\to P\to Q\to R\to0
$$
is called a \textbf{short exact sequence}. Exact sequences
infinite in both directions are called \textbf{long exact
sequences}.
\end{defn}

\begin{prop}\label{sucex}\leavevmode
\begin{enumerate}
\item $0\to A\xrightarrow{\alpha}B$ is exact if and only if
$\ker\alpha=0$, that is, $\alpha$ is injective.
\item $A\xrightarrow{\alpha}B\to0$ is exact if and only if
$\text{img }\alpha=B$, that is, $\alpha$ is surjective.
\item $0\to A\xrightarrow{\alpha}B\to0$ is exact if and only
if $\alpha$ is an isomorphism, by the two preceding items.
\item $0\to A\xrightarrow{\alpha}B\xrightarrow{\beta}C\to0$
is exact if and only if $\alpha$ is injective, $\beta$ is
surjective, and $\ker\beta=\text{img }\alpha$, so that
$\beta$ induces an isomorphism $C\cong B/\text{img }\alpha$.

If we think of $\alpha$ as the inclusion of $A$ as a subgroup
of $B$, we can write $C\cong B/A$.
\end{enumerate}
\end{prop}

\begin{obs}[First isomorphism theorem]
If $M'\subseteq M$, then
$$
0\to M'\to M\to M/M'\to0
$$
is an exact sequence.
\end{obs}

\section{Homotopy}

\begin{defn}
Two homomorphisms
$$
f,g:(C_\bullet,\partial)\to(C'_\bullet,\partial')
$$
are \textbf{homotopic} if there exist homomorphisms
$H_p:C_p\to C'_{p+1}$ for every $p\in\Z$ such that
$$
f_p-g_p=\partial'_{p+1}H_p+H_{p-1}\partial_p.
$$
The omitted diagram visualizes these arrows. Thus the sum of
the blue arrows is equal to the red arrow. We are not saying
that the diagram is commutative.
\end{defn}

\begin{lema}
With the notation above,
$\bar f_p=\bar g_p:H_p(C_\bullet)\to H_p(C'_\bullet)$. That
is, homotopic functions induce equal functions in homology.
\end{lema}

\section{The snake lemma}

\begin{lema}[Snake lemma]
Consider the commutative diagram of $R$-modules and suppose
that its rows are exact. Then there exists a homomorphism
$\delta_*:\ker\partial_3\to Z_1/\img\partial_1$ such that the
usual snake-lemma sequence
$$
\ker\partial_1\to\ker\partial_2\to\ker\partial_3
\xrightarrow{\delta_*}
Z_1/\img\partial_1\to Z_2/\img\partial_2
\to Z_3/\img\partial_3
$$
is exact, where the first maps are restrictions and the last
maps are induced maps.
\end{lema}

\begin{obs}
Intuitively, the cokernel gives us information about how far
a homomorphism is from being surjective.
\end{obs}

\section{The fundamental theorem of homological algebra}

First we introduce some notation.

\begin{defn}
We say that a sequence of chain complexes
$$
\cdots\to C_\bullet\xrightarrow fD_\bullet
\xrightarrow gE_\bullet\to\cdots
$$
is exact at $D_\bullet$ if
$$
\cdots\to C_p\xrightarrow {f_p}D_p
\xrightarrow {g_p}E_p\to\cdots
$$
is exact for every $p\in\Z$.
\end{defn}

\begin{teo}[Fundamental theorem of homological algebra]
If
$$
\cdots\to A_\bullet\xrightarrow\phi B_\bullet
\xrightarrow\psi C_\bullet\to\cdots
$$
is an exact sequence of chain complexes, then there exist
homomorphisms
$$
\delta_{*p}:H_p(C_\bullet)\to H_{p-1}(A_\bullet)
$$
such that the sequence
$$
\cdots\to H_p(A_\bullet)\to H_p(B_\bullet)
\to H_p(C_\bullet)\xrightarrow{\delta_{*p}}
H_{p-1}(A_\bullet)\to H_{p-1}(B_\bullet)\to\cdots
$$
is exact.
\end{teo}

In the following commutative diagram one can see clearly what
is happening. We explain a little how to define the connecting
homomorphism by diagram chasing. Start with a cycle
$c\in C_p(A)$. Since $j_p$ is surjective, there exists an
$a\in B_p$ such that $j_p(a)=c$. Then
$\partial_p(a)\in\ker j_{p-1}$, since the diagram commutes
and $c$ is a cycle. Since the sequence is exact,
$\ker j_{p-1}=\img i_{p-1}$, so there exists
$a\in A_{p-1}$ such that $i_{p-1}(a)=\partial_p(b)$. This
$a$ is a cycle, since the diagram commutes,
$i_{p-2}(a)=\partial(\partial(b))=0$, and $i_{p-2}$ is
injective by exactness. Thus we define $\delta_{*p}[c]=[a]$.

\section{Naturality of the connecting homomorphism}

\begin{teo}[Naturality of the connecting homomorphism]
Given a commutative diagram of short exact sequences of chain
complexes
$$
0\to A_\bullet\to B_\bullet\to C_\bullet\to0,
\qquad
0\to A'_\bullet\to B'_\bullet\to C'_\bullet\to0,
$$
with vertical maps $f,g,h$, the corresponding diagram of long
exact homology sequences commutes.
\end{teo}

It seems that this is a property related to the functorial
structure of homology.

\section{The five lemma}

\begin{lema}[Five lemma]
Consider the commutative diagram with exact rows
$$
M_5\to M_4\to M_3\to M_2\to M_1,
\qquad
N_5\to N_4\to N_3\to N_2\to N_1.
$$
If $h_5,h_4,h_2$ and $h_1$ are isomorphisms, then $h_3$ is
also an isomorphism.
\end{lema}

Where will this be used?

\chapter{Singular homology}

\section{Simplices}

We begin by defining several new concepts. Fix an integer
$n\geq0$. An \textbf{$n$-simplex} is the smallest convex set
in $\R^m$ ($m>n$) containing $n+1$ points
$v_0,\ldots,v_n$ which do not lie in a hyperplane of
dimension less than $n$.

% Pictures of simplices omitted.

We denote it by $[v_0,\ldots,v_n]$ and say that
$v_0,\ldots,v_n$ are its \textbf{vertices}. We can write it
as
$$
[v_0,\ldots,v_n]
=
\{t_0v_0+\cdots+t_nv_n
\mid t_i\geq0,\ t_0+\cdots+t_n=1\}.
$$

The \textbf{standard $n$-simplex} is
$\Delta^n:=[e_1,\ldots,e_n]$, where $e_1,\ldots,e_n$ is the
standard basis of $\R^{n+1}$.

We observe that
$$
\Delta^n=
\{(t_0,\ldots,t_n)\in\R^{n+1}
\mid t_0+\cdots+t_n=1\}.
$$
For us the order of the vertices in $[v_0,\ldots,v_n]$ is
important and should always be kept in mind.

Given an $n$-simplex, we always have the function
$$
\begin{aligned}
(v_0,\ldots,v_n):\Delta^n&\to[v_0,\ldots,v_n]\\
(t_0+\cdots+t_n)&\mapsto t_0v_0+\cdots+t_nv_n.
\end{aligned}
$$
We say that $(t_0,\ldots,t_n)$ are the
\textbf{barycentric coordinates} of the point
$t_0v_0+\cdots+t_nv_n\in[v_0,\ldots,v_n]$.

A \textbf{face} of $[v_0,\ldots,v_n]$ is the subsimplex
generated by any nonempty subset of $\{v_0,\ldots,v_n\}$. We
will regard any 1-dimensional face $[v_i,v_j]$ with $i<j$ as
oriented in increasing order.

How are the 2-dimensional faces oriented?

\section{The singular chain complex}

Take a topological space $X$ and an associative ring with
unit $R$. A \textbf{singular $n$-simplex} is a function
$\sigma:\Delta^n\to X$.

The word ``singular'' comes from the fact that no conditions
are imposed on $\sigma$ except continuity. This means that a
singular simplex may look quite different from what we first
imagine.

Define
$$
C_n(X):=\left\{
\sum_{i=1}^m r_i\sigma_i
\mid m\in\Z,
 r_i\in R,
 \sigma_i\text{ is a singular simplex}
\right\}.
$$
This is the free $R$-module generated by the set of singular
$n$-simplices. The elements of $C_n$ are called singular
$n$-chains. We want to construct the sequence
\begin{equation}\label{eq2.1}
\cdots\to C_{n+1}(X)\xrightarrow{\partial_{n+1}}
C_n(X)\xrightarrow{\partial_n}C_{n-1}(X)
\to\cdots\xrightarrow{\partial_1}C_0(X).
\end{equation}
To do this it is enough to define
$\partial_n:C_n(X)\to C_{n-1}(X)$ as follows: for a singular
$n$-simplex $\sigma:\Delta^n=[v_0,\ldots,v_n]\to X$,
$$
\partial_n(\sigma)=
\sum_{i=0}^n(-1)^i
\sigma|[v_0,\ldots,\hat v_i,\ldots,v_n].
$$
Here $\sigma|[v_0,\ldots,\hat v_i,\ldots,v_n]$ is the
following singular $(n-1)$-simplex: first take the
$(n-1)$-face of $\Delta^n$ obtained by removing the vertex
$v_i$, namely
$$
[v_0,\ldots,\hat v_i,\ldots,v_n]
:=
[v_0,\ldots,v_{i-1},v_{i+1},\ldots,v_n],
$$
and then simply compose with $\sigma$. The bottom arrow in
the omitted diagram is the obvious map: it sends the vertices
in order and skips the $i$-th one. Thus $\partial_n$ is
defined on the basis of $C_n$, and we extend linearly to all
of $C_n$.

Now we see a proposition:

\begin{prop}
$\partial_{n-1}\circ\partial_n=0$.
\end{prop}

Thus the sequence \eqref{eq2.1} is a chain complex, which we
call the \textbf{singular chain complex of $X$}, denoted
$C_\bullet(X)$. We can now consider its homology groups and
define
$$
H_n(X;R):=H_n(C_\bullet(X))
$$
as the \textbf{$n$-th singular homology group of $X$ with
coefficients in $R$}.

\section{First properties of homology}

\section{Homology and path-connected components}

\begin{prop}
Let $X=\bigsqcup X_i$ be the decomposition of the topological
space $X$ into path-connected components. Then
$$
H_n\left(\bigsqcup X_i,R\right)
\cong
\bigoplus H_n(X_i,R).
$$
\end{prop}

\section{The zeroth homology group}

\begin{prop}
For any space $X$, $H_0(X;R)$ is a direct sum of copies of
$R$, one for each path-connected component.
\end{prop}

\section{The homology of a point}

\begin{prop}
If $X$ consists of a single point, then
$$
H_n(X;R)=
\begin{cases}
R\qquad\text{if } n=0,\\
0\qquad\text{if } n\neq0.
\end{cases}
$$
\end{prop}

For this to be fully intuitive: we would like the homology of
this space to be zero in all dimensions. We will fix this in
the next section.

\section{Reduced homology}\label{sec:6.4}

Reduced homology is a variant of singular homology useful for
computations in low dimensions.

\begin{defn}
Let $X$ be a topological space and $R$ an associative ring
with unit. The \textbf{reduced homology} groups of $X$ with
coefficients in $R$, denoted $\tilde H(X;R)$, are the
homology groups of the chain complex
$$
\cdots\to C_p(X)\xrightarrow{\partial_p}C_{p-1}(X)
\to\cdots\to C_1(X)\xrightarrow{\partial_1}C_0(X)
\xrightarrow{\varepsilon}R\to0,
$$
where
$\varepsilon(\sum n_i\sigma_i)=\sum n_i$ is the
\textbf{augmentation map}.
\end{defn}

\begin{obs}
Since the chain complex is identical to the singular chain
complex except at the rightmost terms, reduced homology can
only differ from singular homology in dimension 0. That is,
$\tilde H_n(X;R)=H_n(X;R)$ for every $n\geq1$.
\end{obs}

\begin{ejer}[Reduced homology of a one-point space]
If $X=\{x\}$, then $\tilde H_0(X;R)=0$.
\end{ejer}

\begin{proof}
Recall how we computed the zeroth singular homology of $X$.
Since the last homomorphism in the singular chain complex is
$X\xrightarrow{\partial_0}0$, the kernel is the whole group,
so $H_0(X)=C_0(X)/\img(\partial_1)$. We showed using the
augmentation map that this quotient is isomorphic to $R$.
Thus
$$
0\to\img\partial_1\to C_0(X)\xrightarrow{\varepsilon}R\to0
$$
is a short exact sequence. Hence the reduced chain complex is
exact in degree 0, that is,
$\ker(\varepsilon)=\img(\partial_1)=0$, so
$\tilde H_0(X;R)=0$.
\end{proof}

\begin{ejer}
In general, $H_0(X;R)=\tilde H_0(X;R)\oplus R$.
\end{ejer}

\begin{proof}
By the definition of $\varepsilon$, we know that
$\varepsilon|_{\img\partial_1}=0$.

Recall that, in general, a map out of a group descends to the
quotient by a normal subgroup if its kernel contains that
subgroup. This means that we have an induced map
$$
\bar\varepsilon:H_0(X;R)=C_0(X)/\img\partial_1\to R.
$$
This map is surjective, and its kernel is
$\ker\varepsilon/\img\partial_1=\tilde H_0$. In the language
we have developed, this is the same as saying that we have a
short exact sequence
$$
0\to\tilde H_0(X;R)\to H_0(X;R)
\xrightarrow{\bar\varepsilon}R\to0.
$$
At the end of the day,
$H_0(X;R)/\tilde H_0(X;R)\cong R$. In fact this implies that
$H_0(X;R)\cong\tilde H_0(X;R)\oplus R$. To see this, just
notice that $\bar\varepsilon$ sends an element
$g\in H_0(X;R)$ either to $0$, in which case it lies in
$\tilde H_0(X;R)$, or to an element of $R$, which is the rest
of its image. These two groups have trivial intersection and
generate the direct sum $\tilde H_0(X;R)\oplus R$.
\end{proof}

\section{Functoriality}

In this section we will prove that, in categorical language,
the assignment sending each topological space to its chain
complex is functorial. Also, the assignment sending each
space to its $n$-th homology group is functorial.

In other words, what we will see is that maps between
topological spaces induce maps between singular chain
complexes and therefore also maps between singular homology
groups; moreover, compositions go to compositions and the
identity goes to the identity.

\section{Homotopy invariance}

First recall Definition \ref{def:func-homot}.

\begin{defn}
Two functions $f,g:X\to Y$ are \textbf{homotopic} if there
exists a function $H:X\times I\to Y$ such that
$$
H(x,0)=f(x)\qquad\forall x\in X,
\qquad
H(x,1)=g(x)\qquad\forall x\in X,
$$
and this is denoted by $f\simeq g$.
\end{defn}

\begin{defn}
Two spaces $X$ and $Y$ are \textbf{homotopy equivalent} if
there exist functions called \textbf{homotopy equivalences}
$f:X\to Y$ and $g:Y\to X$ such that
$g\circ f\simeq Id_X$ and $f\circ g\simeq Id_Y$. This is
denoted by $X\simeq Y$.
\end{defn}

Now:

\begin{teo}
If two functions $f,g:X\to Y$ are homotopic, then they induce
the same homomorphism on the $n$-th homology group,
$f_*=g_*:H_n(X)\to H_n(Y)$, for every $n$.
\end{teo}

\begin{coro}
If $f:X\to Y$ is a homotopy equivalence, then the induced map
$f_*:H_n(X)\to H_n(Y)$ is an isomorphism for every $n$.
\end{coro}

\begin{proof}
Using the theorem and the functoriality properties.
\end{proof}

\begin{ejem}
If $X\simeq\{x_0\}$, that is, if $X$ is
\textbf{contractible}, then
$$
H_n(X,R)=
\begin{cases}
R\qquad\text{if }n=0,\\
0\qquad\text{if }n\neq0.
\end{cases}
$$
\end{ejem}


\section{Relative homology}

\section{Construction}

Relative homology is a kind of singular homology which
generalizes reduced homology. The point will be to define the
homology of the space while ``ignoring what happens inside
the subspace.''

Take topological spaces $A\subseteq X$. We say that $(X,A)$
is a \textbf{pair}. Notice that $C_\bullet(A)$ is a
subcomplex of $C_\bullet(X)$, so we can define the relative
complex
$$
C_\bullet(X,A)=C_\bullet(X)/C_\bullet(A).
$$
This simply means that for every $n$,
$$
C_n(X,A)=C_n(X)/C_n(A),
$$
so chains in $A$ become trivial.

It is clear that the $n$-th boundary operator restricted to
$C_n(A)$ maps to $C_{n-1}(A)$, since the boundary of a chain
in $A$ could not leave $A$. This means that the boundary map
is well defined on the quotient, giving a chain complex
$$
\cdots\to C_{n+1}(X,A)\xrightarrow{\partial_{n+1}}
C_n(X,A)\xrightarrow{\partial_n}C_{n-1}(X,A)\to\cdots.
$$
This induces the homology
$$
H_n(X,A)=\ker\partial_n/\text{img }\partial_{n-1}.
$$

The first thing that happens is that we have a short exact
sequence to which we will apply the fundamental theorem of
homological algebra:
$$
0\to C_\bullet(A)\xrightarrow{i}C_\bullet(X)
\xrightarrow{j}C_\bullet(X,A)\to0.
$$
Thus we obtain the \textbf{long exact sequence of the pair}
$$
\cdots\to H_n(A)\xrightarrow{i_{*n}}H_n(X)
\xrightarrow{j_{*n}}H_n(X,A)\xrightarrow{\delta_n}
H_{n-1}(A)\xrightarrow{i_{*n-1}}H_{n-1}(X)\to\cdots.
$$

\begin{obs}
If the homology groups of the pair $C_p(X,A)$ were trivial,
we would automatically have that the map induced by the
inclusion is an isomorphism. In fact, this is an if and only
if. Thus the homology groups of the pair measure how
different the homology groups of $A$ are from those of $X$.
\end{obs}

\begin{obs}
We add that although the map $\delta$ used to complete the
long exact sequence of the pair comes from the fundamental
theorem of homological algebra, when we recall the proof of
that theorem we see that this map acts exactly like the
original boundary operator of $X$.
\end{obs}

\begin{obs}
It is also possible to define the \textbf{reduced homology of
the pair}; see Hatcher. In fact it turns out to be equal to
the homology of the pair just defined, that is,
$\tilde H_n(X,A)=H_n(X,A)$.
\end{obs}

\begin{ejer}
$H_n(X,\{x_0\})=H_n(X)$.
\begin{proof}[Solution]
Using the first and third observations above, we know that
there exists the long exact sequence of the pair in reduced
homology. Since $\tilde H_n(\{x_0\})\cong0$ for every $n$, we
obtain the desired isomorphism.

\textcolor{red}{This proof is in Hatcher, p. 118. In the
video this exercise appears before defining the long exact
sequence of the pair and without the comment on reduced
homology. Is it possible to prove it without these two
things?}
\end{proof}
\end{ejer}

\section{Induced functions}

\begin{obs}
A function between pairs $f:(X,A)\to(Y,B)$ induces a function
$f_\#:C_\bullet(X,A)\to C_\bullet(Y,B)$ and therefore also a
function $f_*:H_n(X,A)\to H_n(Y,B)$.
\end{obs}

\begin{prop}
If two functions between pairs $f,g:(X,A)\to(Y,B)$ are
homotopic through functions between pairs, then the induced
maps $f_*$ and $g_*$ are equal.
\end{prop}

\section{Splitting lemma}\label{subsec:splitting}

Take the case where there is a retraction $r:X\to A$, that
is, $ri=id_A$, where $i$ is the inclusion. This implies that
$(ri)_*=r_*i_*=id_{H_n(A)}$, so $r_*i_*$ is injective, and in
fact we can break the exact sequence of the pair into short
exact sequences of the form
$$
0\to H_n(A)\xrightarrow{i_*}H_n(X)
\xrightarrow{j_*}H_n(X,A)\to0.
$$
But there is more.

\begin{lema}[Splitting lemma]
For any short exact sequence of abelian groups
$$
0\to A\xrightarrow{i}B\xrightarrow{j}C\to0,
$$
the following statements are equivalent:
\begin{enumerate}
\item[(a)] there exists a homomorphism $p:B\to A$ such that
$pi=id_A$;
\item[(b)] there exists a homomorphism $q:C\to B$ such that
$qj=id_C$;
\item[(c)] there exists an isomorphism $B\approx A\oplus C$
which makes the corresponding diagram commute, with the
obvious maps in the last row, $a\mapsto(a,0)$ and
$(a,c)\mapsto c$.
\end{enumerate}
In this case we say that the sequence \textbf{splits}.
\end{lema}

\begin{obs}\leavevmode
\begin{itemize}
\item $C$ is free, even if it is not abelian, if and only if
the sequence splits.
\item A retraction $r:X\to A$ gives
$H_n(X)\approx H_n(A)\oplus H_n(X,A)$.
\end{itemize}
\end{obs}

\section{The long exact sequence of a triple}

If $Z\subseteq Y\subseteq X$ are topological spaces, it is
immediate from the definition that
$$
C_\bullet(Z)\leq C_\bullet(Y)\leq C_\bullet(X).
$$
Taking the quotient by $Z$, we obtain
$$
C_\bullet(Y,Z)\hookrightarrow C_\bullet(X,Z)
\to
\frac{C_\bullet(X,Z)}{C_\bullet(Y,Z)}=C_\bullet(X,Y).
$$
In the last equality we use an isomorphism theorem to see
that
$$
\frac{C_\bullet(X)/C_\bullet(Z)}
{C_\bullet(Y)/C_\bullet(Z)}
=
C_\bullet(X,Y).
$$
Applying the fundamental theorem of homological algebra gives
the \textbf{long exact sequence of the triple}
$$
\cdots\to H_n(Y,Z)\xrightarrow{i_*}H_n(X,Z)
\xrightarrow{j_*}H_n(X,Y)\xrightarrow{\delta}
H_{n-1}(Y,Z)\xrightarrow{i_*}\cdots.
$$

\section{Excision}\label{sec:escisión}

This result makes much more precise the idea that homology of
a pair is the homology of the space while ``ignoring what
happens in the subspace.''

Let $Z\subseteq A\subseteq X$ be such that the closure of $Z$
is contained in the interior of $A$. Then
$$
(X-Z,A-Z)\hookrightarrow (X,A)
$$
induces an isomorphism
$$
H_n(X-Z,A-Z)\to H_n(X,A)
\qquad\forall n.
$$
Equivalently, for subspaces $A,B\subseteq X$ whose interiors
cover $X$, the inclusion
$$
(B,A\cap B)\hookrightarrow (X,A)
$$
induces an isomorphism
$$
H_n(B,A\cap B)\to H_n(X,A)
\qquad\forall n.
$$

\section{The homology of a quotient}

Given a topological space $X$ and a subspace $A\subseteq X$,
we can consider the space $X/A$, obtained by identifying all
points of $A$. This space is like $X$ but ignoring everything
that happens in $A$. This reminds us of homology of a pair,
and the natural question arises: is there any relation
between $H_n(X,A)$ and $H_n(X/A)$? We need some conditions.

\begin{defn}
We say that $(X,A)$ is a \textbf{good pair} if $A$ is closed
and is a strong deformation retract of some neighbourhood in
$X$.
\end{defn}

The Hawaiian earring with the attaching point is not a good
pair because every neighbourhood of that point contains an
entire circle, so it cannot be deformation retracted.

\begin{teo}
Let $(X,A)$ be a good pair. Then
$$
q:(X,A)\to(X/A,A/A)
$$
induces isomorphisms for every $n$ of the form
$$
q_*:H_n(X,A)\to H_n(X/A,A/A)
\cong\tilde H_n(X/A),
$$
where the tilde denotes reduced homology.
\end{teo}

The proof uses the five lemma, homotopy invariance, etc.

\begin{coro}
Let $(X,A)$ be a good pair. Then we have the following long
exact sequence:
$$
\cdots\to\tilde H_n(A)\to\tilde H_n(X)\to\tilde H_n(X/A)
\to\tilde H_{n-1}(A)\to\tilde H_{n-1}(X)
\to\tilde H_{n-1}(X/A)\to\cdots.
$$
\end{coro}

\section{Homology of the sphere and applications}

\begin{teo}[Homology of the sphere]
$$
\tilde H_i(S^n;R)=
\begin{cases}
R\qquad\text{if }n=i,\\
0\qquad\text{if }n\neq i.
\end{cases}
$$
\end{teo}

\begin{teo}[Brouwer fixed point theorem]
Let $n\geq2$ and $f:D^n\to D^n$. Then $f$ has a fixed point.
\end{teo}

\begin{teo}[Homology of the wedge]
Consider a wedge of spaces $\bigvee_{\alpha\in I}X_\alpha$
such that $(X_\alpha,x_\alpha)$ is a good pair for every
$\alpha$, where $x_\alpha$ is the attaching point. Then the
inclusions $i_\alpha:X_\alpha\hookrightarrow
\bigvee_{\alpha\in I}$ induce an isomorphism
$$
\bigoplus i_{\alpha *}:
\bigoplus_{\alpha\in I}\tilde H_n(X_\alpha)
\to
\tilde H_n\left(\bigvee_{\alpha\in I}\right).
$$
\end{teo}

\begin{teo}[Invariance of dimension]
Let $U\subseteq\R^n$ and $V\subseteq\R^m$ be open sets. If
$U$ and $V$ are homeomorphic, then $n=m$.
\end{teo}

\section{The Mayer--Vietoris sequence}

Let $X$ be a topological space and let $A,B\subseteq X$ be
such that $X=\Int A\cup\Int B$. Then we have the following
long exact sequence in reduced homology:
$$
\cdots\to H_n(A\cap B)\xrightarrow{\Phi}
H_n(A)\oplus H_n(B)\xrightarrow{\Psi}H_n(X)
\xrightarrow{\delta}H_{n-1}(A\cap B)\to\cdots.
$$

To understand where this sequence came from, we should look
at the chain subgroups $C_n(A+B)$ formed by sums of chains in
$A$ and chains in $B$. It is possible to construct a chain
complex, since the boundary homomorphism of the whole space,
$\partial:C_n(X)\to C_{n-1}(X)$, restricts correctly to
$C_n(A+B)\to C_{n-1}(A+B)$. Using the machinery of the
excision theorem, one proves that the inclusion
$C_n(A+B)\hookrightarrow C_n(X)$ induces isomorphisms on
homology. Then, for every $n$, we have the short exact
sequence
$$
0\to C_n(A\cap B)\xrightarrow{\varphi}
C_n(A)\oplus C_n(B)\xrightarrow{\psi}C_n(A+B)\to0,
$$
where $\varphi(x)=(x,-x)$ and $\psi(x,y)=x+y$.

This way of defining the functions ensures that the sequence
is exact. It is easy to see that $\psi\varphi=0$, so
$\img\varphi\subseteq\ker\psi$. Also
$\ker\psi\subseteq\img\varphi$: if
$(x,y)\overset{\psi}{\mapsto}(0,0)$, then $x=-y$, so both
$x$ and $y$ lie in $A$ and in $B$, that is,
$(x,y)\in C_n(A\cap B)$; clearly it lies in the image of
$\varphi$. It is also easy to see that $\ker\varphi=0$ and
that $\img\psi=C_n(A+B)$.

Applying the fundamental theorem of homological algebra to
the short exact sequence of chain complexes induced by these
short exact sequences in each degree gives the Mayer--Vietoris
sequence.

It is possible to say how the augmentation map acts. By what
we said about excision, an element
$[w]\in H_n(X)\cong H_n(A+B)$ is of the form $[a+b]$ for
chains $a\in C_n(A)$ and $b\in C_n(B)$ whose sum $a+b$ is a
cycle, that is, $\partial(a+b)=0$. Then
$\partial(a)=-\partial(b)$ is a cycle in $H_{n-1}(A\cap B)$,
so it is the perfect representative for the equivalence class
to which we send $[w]$.

Everything said so far is summarized by
$$
\begin{aligned}
i:A\cap B&\hookrightarrow A,
&j:A\cap B&\hookrightarrow B,\\
\Phi(x)&=(i_*x,-j_*x),
&\Psi(y,z)&=k_*y+\ell_*z,\\
[w]=[a+b]&\overset{\delta}{\mapsto}[
\partial(a)]=[-\partial(b)]
& &\text{for }a\in C_n(A),\ b\in C_n(B).
\end{aligned}
$$

\begin{ejem}
It is possible to compute the homology of the torus using the
Mayer--Vietoris sequence, although it is not enough to write
down the exact sequence: one has to understand how the
homomorphisms act.
\end{ejem}

\begin{obs}
There is a version of this sequence for $X=A\cup B$ where
$A$ and $B$ are deformation retracts of certain
neighbourhoods $U$ and $V$ and where $U\cap V$ retracts onto
$A\cap B$.

In particular, if $X$ is a CW complex and $A,B\subseteq X$
are CW subcomplexes, then it is possible to prove that such
neighbourhoods exist, and we obtain the same Mayer--Vietoris
sequence
$$
\cdots\to\tilde H_n(A\cap B)\to
\tilde H_n(A)\oplus\tilde H_n(B)\to\tilde H_n(X)
\to\tilde H_{n-1}(A\cap B)\to\cdots.
$$
Now $A$ and $B$ do not have to be open, and it does not have
to be true that $\Int A\cup\Int B=X$.
\end{obs}

\begin{ejem}\label{ejem:toroCWMV}
\textcolor{red}{In my class notes there is a way to compute
$H_1$ and $H_0$ of the torus viewed as a CW complex with
subcomplexes.}

When we view the torus $S^1\times S^1$ as the quotient of a
square with sides identified, we have a CW-complex structure
with one 0-cell, two 1-cells, and one 2-cell. Here, we can
take $A=D^2$ and $B=S^1\vee S^1$.

Almost all homology groups become zero
\textcolor{red}{I do not see why}; here one has to use that
homology of a wedge is the sum of the homologies of the
wedge summands, and that the homology of $S^1$ is zero when
the subscript is greater than or equal to 2.

Anyway, for $m>2$, $\tilde H_m(X)=0$. It is easy to see that
$H_1(S^1\times S^1)\cong R\oplus R$. But for
$H_2(S^1\times S^1)$ something is missing.
\end{ejem}

\section{Representatives of the homology of the sphere}

In this section we use $R=\mathbb Z$.

Our method will be to find generators for the relative
homology $H_n(D^n,\partial D^n)$, since this should be
isomorphic to $H_n(S^n)$. The first thing we do is replace
that pair by the following one, to which it is also
isomorphic: $(\Delta^n,\partial\Delta^n)$.

\begin{prop}
The identity map $Id:\Delta^n\to\Delta^n$ is a cycle which
generates $H_n(\Delta^n,\partial\Delta^n)\cong\Z$.
\begin{proof}
First note that the identity is indeed a cycle in the
relative homology group, since the kernel of the boundary map
$C_n(\Delta^n)\xrightarrow{\partial_n}C_{n-1}(\Delta^n)$ is
contained in the subgroup
$C_n(\partial\Delta^n)\leq C_n(\Delta^n)$. A group
homomorphism descends to a quotient if the kernel is contained
in the normal subgroup.

Now we do induction on $n$. (...)

% If $n=0$, then $\Delta^0=\{e_0\}$ is a point and its
% boundary is empty. Hence
% $H_0(\Delta^0,\partial\Delta^0)=H_0(\Delta^0)=
% \mathbb Z=\langle[e_0]\rangle$.
% Now suppose the statement for $n$ and prove it for $n+1$.
% Define $\Lambda$ as the union of all the $(n-1)$-faces in
% the boundary except one. Then ...
\end{proof}
\end{prop}

Now let us think of $S^n$ as the union of two standard
simplices glued along their boundary, denoted
$\Delta_1^n\cup_{\partial\Delta^n}\Delta_2^n$.

\begin{prop}
$\Delta^n_1-\Delta^n_2$ is a generator of $H_n(S^n)$ for
$n>0$.
\end{prop}

\section{Degree of a map between spheres}

Consider $f:S^n\to S^n$, which induces a function
$$
f_*:\tilde H_n(S^n)\cong\mathbb Z
\to
\tilde H_n(S^n)\cong\mathbb Z,
$$
so that $f_*$ has no choice but to be of the form
$f_*(n\alpha)=dn\alpha$.

\begin{defn}
In the notation above, the \textbf{degree} of $f$ is $d$.
\end{defn}

\begin{prop}\leavevmode
\begin{enumerate}
\item $\deg Id=1$, by functoriality.
\item If $f$ is not surjective, then $\deg(f)=0$. Since $f$
is not surjective, choose a point $x\notin\text{img }f$.
The relevant diagrams show that $f_*$ factors through
$\tilde H_n(S^2-x)\cong0$, so $f_*$ is constant.
\item If $f\simeq g$, then $\deg f=\deg g$. By homotopy
invariance, the induced functions are equal.
\item $\deg(fg)=\deg(f)\deg(g)$.
\item If $f$ is a reflection interchanging the hemispheres,
then $\deg(f)=-1$.

This is because, as seen in the previous section, the
$n$-sphere, viewed as the union of two standard simplices,
$S^n=\Delta_1^n\cup_{\partial\Delta^n}\Delta_2^n$, has
$n$-th homology generated by the chain
$\Delta_1^n-\Delta_2^n$. The reflection interchanging these
simplices sends the generator to its negative.
\item The antipodal map $-Id:S^n\to S^n$ has degree
$(-1)^{n+1}$. That is,
$$
\deg(-Id)=
\begin{cases}
1\qquad\text{if }n\text{ is odd},\\
-1\qquad\text{if }n\text{ is even}.
\end{cases}
$$
\item If $f$ has no fixed points, then
$\deg(f)=(-1)^{n+1}$. Moreover, if $f$ has no fixed points,
it is homotopic to the identity map.

If $x\neq f(x)$, the line segment in $\R^{n+1}$ passing
through $f(x)$ and $-x$, given by $(1-tf(x))-tx$, does not
pass through the origin. This allows us to define a homotopy
$(1-tf(x))-tx/|(1-tf(x))-tx|$ between $f$ and the antipodal
map.
\end{enumerate}
\end{prop}

\section{Local degree}\label{subsec:grloc}

Following Hatcher:

\begin{defn}
Suppose that $f:S^n\to S^n$ is such that the inverse image of
some $y\in S^n$ consists of finitely many points, say
$f^{-1}(y)=\{x_1,\ldots,x_m\}$. Take a family of disjoint
open sets $U_1,\ldots,U_m$ containing each $x_i$, and which
map to a neighbourhood $V$ of $y$ under $f$. Through a
commutative diagram, using excision and the long exact
sequence, Hatcher argues that the homomorphism $f_*$
originally defined as
$$
H_n(S^n,S^n-x_i)\xleftarrow{\cong}
H_n(U_i,U_i-x_i)\xrightarrow{f_*}H_n(V,V-y)
$$
is of the form $f_*:\Z\to\Z$, multiplication by a number
called the \textbf{local degree} of $f$, denoted
$\deg f|x_i$.
\end{defn}

And then:

\begin{teo}
$\deg f=\sum_i\deg f|_{x_i}$.
\end{teo}

We add another proposition.

\begin{prop}
The degree of the suspension $Sf:S^{n+1}\to S^{n+1}$ of a
function $f:S^n\to S^n$ is equal to the degree of $f$.
\end{prop}

\section{The hairy ball theorem}

\begin{teo}[Hairy ball theorem]
$S^n$ has a vector field which vanishes nowhere if and only
if $n$ is odd.
\end{teo}

\section{Free actions on the sphere}

\begin{teo}[Free actions on the sphere]
$\Z/2\Z$ is the only nontrivial group which acts freely on
$S^n$ if $n$ is even.
\end{teo}

\chapter{CW complexes}\label{chap:CCW}

\section{Construction and basic properties}

Let us establish some notation.
\begin{itemize}
\item $D^n=\{x\in\mathbb R^n:|x|\leq1\}$ is the closed
$n$-disk.
\item $S^n=\{x\in\mathbb R^n:|x|=1\}$ is the $n$-sphere.
\item $e^n=\{x\in\mathbb R^n:|x|<1\}$ is the open $n$-cell,
or simply the $n$-cell.
\end{itemize}
Thus
\begin{itemize}
\item $\partial D^n=S^{n-1}$;
\item $D^n=e^n\cup S^{n-1}$ as sets, not using the disjoint
union topology;
\item and, well, it should be true that $e^0=\{pt\}$.
\end{itemize}

\begin{defn}
A space $X$ is a CW complex if it can be constructed by the
following procedure:
\begin{enumerate}
\item Start with a discrete space $X^0$, called the
0-skeleton.
\item Obtain the $n$-skeleton $X^n$ from the $(n-1)$-skeleton
$X^{n-1}$ by gluing $n$-cells $e^n_\alpha$ via maps
$\varphi_\alpha:S^{n-1}\to X^{n-1}$. Formally,
$$
X^n=X^{n-1}\bigsqcup_\alpha D^n_\alpha\Big/\sim,
\qquad
x\sim\varphi_\alpha(x)
\quad\text{if }x\in\partial D^n_\alpha.
$$
\textit{The point is to take points and join them with lines,
and then take lines and fill them with disks, etc.}
\item We could stop after finitely many steps, say $m$, and
obtain $X=X^m$, but we can also continue gluing cells of
arbitrarily large dimension and obtain $X=\bigcup_nX^n$. In
this case $X$ has the weak topology: $A\subseteq X^n$ is open
(closed) iff $A\cap X^n$ is open (closed) for all $n$.
\end{enumerate}
\end{defn}

\begin{defn}
If $X=X^n$ for some $n\in\mathbb N$, we say that $X$ is
finite-dimensional and define the dimension of $X$ as the
dimension of the largest cell adjoined. Formally,
$\dim X=\min\{n:X^n=X\}$.
\end{defn}

\begin{ejem}\leavevmode
\begin{itemize}
\item The 0-dimensional CW complexes are graphs.
\item $S^n$ is a CW complex.
\item Real projective space. Initially defined as
$\R P^n=\R^{n+1}-\{0\}/x\sim\lambda x$ for $\lambda\neq0$,
we can take representatives of the lines on the sphere to
obtain $\R P^n=S^n/x\sim -x$.

This is equivalent to taking either hemisphere of the sphere,
$D^n$, and identifying the points of its boundary by the
antipodal map. In turn, this is equivalent to taking the
$(n-1)$-sphere, which is the boundary of the disk,
identifying antipodal points, and adding the interior of the
$n$-disk. And the $(n-1)$-sphere with that identification is
precisely $\R P^{n-1}$.

In summary, by induction,
$$
\begin{aligned}
\R P^n
&=\R^{n+1}-\{0\}/x\sim\lambda x,
\qquad\lambda\neq0\\
&=S^n/x\sim-x\\
&=D^n/x\sim-x,
\qquad x\in\partial D^n=S^{n-1}\\
&=\R P^{n-1}\cup e^n\\
&=\R P^{n-2}\cup e^{n-1}\cup e^n\\
&=e^0\cup e^1\cup\ldots\cup e^n,
\end{aligned}
$$
where the attaching maps are exactly the antipodal maps.
\item Complex projective space. As in the real case, complex
projective space, initially defined as
$\C P^n=\C^{n+1}-\{0\}/x\sim\lambda x$ for $\lambda\neq0$,
can be defined only on the sphere $S^{2n+1}$. Just like in
the real case, it is possible to take a disk $D^{2n}$ and
identify antipodal points on its boundary. For now we do not
specify how, but it turns out that
$$
\C P^n=e^0\cup e^2\cup\ldots e^{2n}.
$$
\end{itemize}
\end{ejem}

\begin{ejer}
Every CW complex is semilocally simply connected. That is,
every CW complex has a universal covering.
\end{ejer}

\begin{defn}
For each $n$-cell $e^n_\alpha$ of a CW complex, we have the
\textbf{characteristic map}, which is of the form
$\Phi_\alpha:D^n_\alpha\to X^n$ and extends the attaching map
$\varphi_\alpha$. It is a homeomorphism from the interior of
$D^n_\alpha$ onto $e^n_\alpha$.
\end{defn}

\begin{defn}
Let $X$ be a CW complex. A \textbf{CW subcomplex} $Z$ is a
closed subspace which is also a union of cells of $X$.
\end{defn}

\begin{obs}
For a subcomplex $Z$,
\begin{itemize}
\item $Z\subseteq X$ is a good pair, meaning $Z$ is a
deformation retract of a neighbourhood in $X$;
\item $X/Z$, the space obtained by collapsing $Z$ to a point,
is a CW complex.
\end{itemize}
\end{obs}

\begin{obs}[Very important]
Let $X$ be a CW complex. The quotient by the $(n-1)$-skeleton
is a wedge of spheres, as many as there are cells in the
$n$-skeleton. In symbols,
$$
X^n/X^{n-1}=\bigvee_{\alpha\in I_n}S^n,
$$
where $I_n$ is the set indexing the $n$-cells.
\end{obs}

\begin{teo}
If $X$ and $Y$ are CW complexes with at most countably many
cells, then $X\times Y$ is a CW complex.
\end{teo}

Here, the dimension of the product of two cells, say
$e^1_\alpha\times e^2_\beta$, is the sum of the dimensions of
each one. In this example, we have a 3-cell.

\begin{defn}
Let $X$ be a space. The \textbf{suspension} of $X$ is the
space
$$
SX=X\times I/\sim,
$$
where the equivalence relation $\sim$ identifies the caps,
that is,
$$
(x,1)\sim(y,1),
\qquad
(x,0)\sim(y,0)
\quad\forall x,y\in X.
$$

The \textbf{cone} of $X$ is obtained by identifying only one
cap:
$$
CX=X\times I/\sim,
\qquad
(x,1)\sim(y,1)
\quad\forall x,y\in X.
$$

It is also possible to define the \textbf{suspension of a
function} $f:X\to Y$, which is of the form
$Sf:SX\to SY$, simply by taking the quotient of
$f\times Id:X\times I\to Y\times I$.
\end{defn}

\begin{teo}
The cone and suspension of a CW complex are CW complexes.
\end{teo}

\section{CW complexes according to Lee (extra)}

Let us see how Lee defines CW complexes in his book
\textit{Introduction to Topological Manifolds}.

\begin{defn}
If $X$ is a nonempty topological space, a \textbf{cellular
decomposition} of $X$ is a partition $\mathcal E$ of $X$ into
subspaces which are open cells, such that for each
$e\in\mathcal E$ of dimension $n\geq1$ there exists a
\textbf{characteristic map} $\phi$ from a closed $n$-cell $D$
to $X$, restricting to a homeomorphism from $\text{Int}D$ to
$e$, and sending $\partial D$ into the union of cells in
$\mathcal E$ of dimension strictly less than $n$.

A \textbf{cellular complex} is a Hausdorff space $X$ together
with a cellular decomposition.
\end{defn}

\begin{obs}
Although the cells are homeomorphic to open disks, they do
not have to be open subsets of $X$.
\end{obs}

\begin{defn}
A \textbf{CW complex} is a cellular complex $(X,\mathcal E)$
such that:
\begin{enumerate}
\item[\textbf{C}] \textbf{Closure finiteness.} The closure of
each cell intersects only finitely many cells of $X$.
\item[\textbf{W}] \textbf{Weak topology.} The topology of $X$
is coherent with the family of closed subspaces
$\{\bar e:e\in\mathcal E\}$, that is, $U\subseteq X$ is closed
in $X$ if and only if $U\cap\bar e$ is closed for every
$e\in\mathcal E$.
\end{enumerate}
An equivalence between this definition and Hatcher's, when
$X$ is Hausdorff, is in Hatcher's appendix.
\end{defn}

\section{The singular homology of a CW complex}

\begin{lema}\leavevmode
\begin{itemize}
\item If the $n$-cells are indexed by the set $I_n$, then
$$
H_k(X^n,X^{n-1})=
\begin{cases}
0\qquad\text{if }k\neq n,\\
\bigoplus_{\alpha\in I_n}R\qquad\text{if }k=n.
\end{cases}
$$
\begin{proof}
$H_k=\tilde H_k(X^n/X^{n-1})
=\tilde H_k(\bigvee_{\alpha\in I_n}S^n)
=\bigoplus_{\alpha\in I_n}H_k(S^n)$.
\end{proof}
\item $H_k(X^n)=0$ for every $k>n$.

Again, homology sees only up to the dimension of the space.
\item The inclusion $X^n\hookrightarrow X$ induces an
isomorphism $H_k(X^n)\to H_k(X)$ for $k>n$.

Here the idea is that attaching cells of dimension greater
than the homology we are computing does not change homology.
That is, for $n>k$ one has $H_k(X\cup D^n)=H_k(X)$.
\end{itemize}
\end{lema}

\begin{proof}
Look at $(X^n,X^{n-1})$. There is a short exact sequence of
the pair:
$$
H_{k+1}(X^n,X^{n-1})\to H_k(X^{n-1})
\to H_k(X^n)\to H_k(X^n,X^{n-1}).
$$
Using item 1, if $k\neq n,n-1$, then
$H_k(X^{n-1})\cong H_k(X^n)$.

To prove (2), we have $k>n$, and then, fixing $k$ and going
down one by one, $H_k(X^n)\cong H_k(X^0)$, which is a
discrete space, and homology in positive degree of discrete
spaces is zero. This finishes (2).

For (3), if $k<n<n+1<n+2<\cdots$, then simply
$$
H_k(X^n)\cong H_k(X^{n+1})\cong\cdots\cong H_k(X^{n+m}),
$$
and if $X$ is finite-dimensional, we eventually stop and we
are done. The infinite-dimensional case is left for the
future.
\end{proof}

\section{Cellular homology}

Next we construct a new type of homology for CW complexes.

\begin{defn}
Consider the maps obtained from the long exact sequences of
the pairs of skeleta. The horizontal maps $d_i$ satisfy
$$
d_{n+1}\circ d_n=0,
$$
so we may call the resulting complex the
\textbf{cellular chain complex} $C_\bullet^{CW}(X,R)$. We
then define \textbf{cellular homology} by
$$
H_n^{CW}(X;R)=H_n(C_\bullet^{CW}(X)).
$$
\end{defn}

\begin{teo}
$$
H_n^{CW}(X;R)\cong H_n(X;R)
\qquad\forall n\geq0.
$$
\end{teo}

The first depends on the cellular structure of the space $X$,
but the second does not. Thus the homology of a CW complex is
independent of the cellular structure.

\begin{obs}\label{obs1}
These groups $H_n(X^n,X^{n-1})$ are direct sums of $R$, one
for each $n$-cell of $X$.
\end{obs}

\begin{prop}[Consequences of the definition]
Take $X$ a CW complex and $R=\Z$.
\begin{enumerate}
\item $H_n(X)=0$ if $X$ has no $n$-cells.
\item Suppose $X$ has exactly $k$ $n$-cells. Then $H_n(X)$ is
generated by at most $k$ elements, since
$H_n(X^n,X^{n-1})$ is free abelian on $n$ elements, and so
are the subgroup $\ker d_n$ and the quotient
$\ker d_n/\img d_{n-1}=H_n(X)$.
\item If $X$ has no cells in adjacent dimensions, then
$H_n(X)$ is free for every $n$, with rank equal to the number
of $n$-cells.
\end{enumerate}
\end{prop}

\section{A formula for the boundary homomorphism}

Now let us try to understand what $d$ is. Consider the usual
map from the boundary sphere of an $n$-cell to
$X^{n-1}$, followed by the quotient
$X^{n-1}\to X^{n-1}/X^{n-2}=\bigvee_{\gamma\in I_n}S^{n-1}_\gamma$
and projection to one sphere. The function obtained,
$\varphi_{\alpha\beta}$, is a map between spheres and has a
degree $d_{\alpha\beta}$.

\begin{teo}
For $n\geq2$, the homomorphism $d_n$ sends
$$
e^n_\alpha\mapsto\sum d_{\alpha\beta}e^{n-1}_\beta.
$$
\end{teo}


\section{Examples}

\begin{ejem}
We compute the homology of
$X=\bigvee_{\alpha\in I}S^1_\alpha$ in two different ways:
first using directly that $H_0(X)=R$, and then using the
preceding theorem. Since there is only one 0-cell, for every
$\alpha$, $e^1_\alpha\mapsto x_0-x_0=0$, so in fact
$d_1=0$. It follows that $H_1=\bigoplus_IR$.
\end{ejem}

In fact, this is true in general:

\begin{lema}
If $X$ is a complex with a single 0-cell, then $d_1=0$.
\end{lema}

\begin{ejem}[The torus]\label{ejem:toroCW}
We know that the torus has a CW decomposition with one
0-cell, two 1-cells, and one 2-cell. This gives the chain
complex
$$
0\to\Z\xrightarrow{d_2}\Z\oplus\Z\xrightarrow{d_1=0}\Z\to0.
$$
We know $d_1=0$ because there is only one 0-cell. To discover
how $d_2$ acts, look at the omitted square diagram.

The boundary of the 2-cell, $S^1$, maps by the attaching map
to the wedge of two circles, and \textbf{traverses each one
twice and in opposite directions}. The quotient $X^1/X^0$
does nothing, since $X^0$ was already just one point. In the
last step, when choosing one of the two circles in the wedge,
we realize that our map traverses it \textbf{out and then
back}. This map is homotopic to a constant, so it has degree
zero. Doing this for the two circles, we conclude that
$d_2=0$. Hence $H_1(T^2)=\Z\oplus\Z$ and $H_2(T^2)=\Z$.
\end{ejem}

\begin{teo}[Extra]
A compact connected surface can only be a sphere, a torus, a
connected sum of tori, the projective plane, or a connected
sum of projective planes.
\end{teo}

\begin{ejem}[The torus of genus $g$]
It turns out that the double torus can be seen as the
quotient of an octagon with identified edges. In general, the
$g$-torus is the quotient of a $4g$-gon. The final result is
$$
H_0(S_g)=\Z,
\qquad
H_1(S_g)=\Z^{2g},
\qquad
H_2(S_g)=\Z.
$$
\end{ejem}

\begin{ejem}[The projective plane]
Recall the simplest decomposition of the projective plane,
with one cell in each dimension. Here,
$$
0\to\Z\xrightarrow{d_2}\Z\xrightarrow{d_1=0}\Z\to0.
$$
Formally, to discover how the attaching map acts, we reason
using local degree. The map in the drawing sends the northern
hemisphere, without its boundary, to the circle on the right
without the point. The restriction of the map to the northern
hemisphere helps decompose the function in order to compute
the degree. Since the restriction is a homeomorphism, the
local degree here is 1.

The same happens for the piece of the function on the
southern hemisphere. The key observation is that each of
these local functions is obtained from the other by composing
with the antipodal map, which has degree $1$ on $S^1$. The
degree of the whole function is the sum of the local degrees,
so the degree is 2. Since there is only one cell in each
dimension, $d_2=2$. Hence
$$
H_2(\R P^2)=0,
\qquad
H_1(\R P^2)=\Z/2.
$$
\end{ejem}

\begin{ejem}[$\R P^n$]
We only have to generalize the preceding example. The CW
decomposition of $\R P^n$ is
$$
\R P^n=e_0\cup e_1\cup\cdots\cup e_n.
$$
The diagram above becomes the corresponding diagram involving
$S^{k-1}$, $\mathbb RP^{k-1}$, and the quotient
$\mathbb RP^{k-1}/\mathbb RP^{k-2}$.

The attaching maps $\varphi_\alpha$ are the quotient
projection by the antipodal map on the sphere, i.e. the
covering of the projective plane. These maps are two-to-one,
and this makes the northern hemisphere of $S^{k-1}$ map
homeomorphically to $S^{k-1}$ minus one point, just as in
dimension 2.

As we said, the functions obtained by restricting
$\bar\varphi$ to each hemisphere differ from each other by the
antipodal map. When computing degree, we obtain the formula
$$
\deg(\bar\varphi)=1+(-1)^k.
$$
Therefore
$$
d_k=\begin{cases}
2\quad\text{if }k\text{ is even},\\
0\quad\text{if }k\text{ is odd}.
\end{cases}
$$
This gives the corresponding cellular chain complexes for
$n$ odd and $n$ even. Thus
$$
H_0(\R P^n)=\Z,
$$
$H_k(\R P^n)=\Z/2$ for odd $k<n$, and $H_n(\R P^n)$ is $0$
if $n$ is even and $\Z$ if $n$ is odd.
\end{ejem}

\begin{ejem}[First example where things change if we change
rings]
Take $\R P^n$ and the ring $R=\Q$. Now the map
multiplication by 2 sends $\Q$ onto all of $\Q$, something
that did not happen with $\Z$. This makes the homology
trivial, that is,
$$
H_n(\R P^n,\Q)=0
\qquad\text{for every } n\geq1.
$$
\end{ejem}

\begin{ejem}
What if we take $R=\Z/2$? Multiplication by 2 is the zero map,
so
$$
H_m(\R P^n,\Z/2)=\Z/2
\qquad 0\leq m\leq n.
$$
\end{ejem}

\begin{ejem}[The M\"obius band]
Take the cellular decomposition of the M\"obius band from the
omitted drawing. It corresponds to the chain complex
$$
0\to\Z\xrightarrow{d_2}\Z^3\xrightarrow{d_1}\Z^2\to0.
$$
In this case $d_1$ is not automatically zero because we have
two 0-cells. However, we know that the band is connected, so
$H_0\approx\Z$.

We can compute $d_1$ explicitly, because we know that it
behaves exactly like the boundary in simplicial homology.
According to the names in the diagram, we can simply see that
$$
\begin{array}{c@{\hspace{1cm}}c@{\hspace{1cm}}c}
d_1(e^1_1)=e^0_2-e^0_1&
 d_1(e^1_2)=e^0_1-e^0_2&
 d_1(e^1_3)=e^0_2-e^0_1.
\end{array}
$$
Equivalently,
$$
\begin{array}{c@{\hspace{1cm}}c@{\hspace{1cm}}c}
(1,0,0)\mapsto(-1,1)&
(0,1,0)\mapsto(1,-1)&
(0,0,1)\mapsto(-1,1).
\end{array}
$$
Thus $\img d_1\approx\Z$, as expected, and
$$
\ker d_1\approx\langle(1,1,0),(0,1,1)\rangle
\approx\Z\oplus\Z.
$$

To compute $d_2$, notice that the 2-cell is attached only
along the edges $e^1_2$ and $e^1_3$. Collapsing the
0-skeleton, we get a wedge of three circles, one for each
edge. The functions induced in our commutative diagram have
degree 0 for $e^1_1$ and degree 1 for $e^1_2$ and $e^1_3$.
Thus
$$
d_2(e^2)=0\cdot e^1_1+e^1_2+e^1_3,
$$
or $1\overset{d_2}{\mapsto}(0,1,1)$. Hence
$\img d_2\approx\Z$ and $\ker d_2\approx0$. Therefore
$$
H_1(M)\approx(\Z\oplus\Z)/\Z\approx\Z,
\qquad
H_2(M)\approx0.
$$
\end{ejem}


\section{Euler characteristic}

Let $Y$ be a finite CW complex, i.e. one with finitely many
cells. In fact this property is equivalent to $Y$ being
compact. We define the Euler characteristic by
$$
\begin{aligned}
\chi(Y)&=(\#\text{ 0-cells})-(\#\text{ 1-cells})
 +(\#\text{ 2-cells})-\cdots\\
&=\sum_{i=0}^{\infty}(-1)^i
(\#\text{ }i\text{-cells})\in\Z.
\end{aligned}
$$

\begin{ejems}\leavevmode
\begin{itemize}
\item
$\chi(S^n)=\begin{cases}
0\qquad n\text{ odd},\\
2\qquad n\text{ even};
\end{cases}$
\item
$\chi(S^n)=\begin{cases}
0\qquad n\text{ odd},\\
2\qquad n\text{ even};
\end{cases}$
\item $\chi(\bigvee_{i=0}^nS^1)=1-n$;
\item $\chi(S_g)=1-2g+1=2-2g$.
\end{itemize}
\end{ejems}

\begin{teo}
Let $X,Y$ be finite CW complexes. If $X\simeq Y$, then
$\chi(X)=\chi(Y)$.
\end{teo}

\begin{prop}
Let $Y$ be a finite CW complex. Then
$$
\chi(Y)=\sum_{i=0}^{\infty}(-1)^i\ran H_i(Y;\Z).
$$
To understand what $\ran H_i$ is, note that the chain complex
is made of finitely generated abelian groups, since $Y$ is
finite. In fact, the homology of $Y$ is an exact sequence of
finite finitely generated groups. Then, by the classification
theorem for finitely generated abelian groups, we can
decompose a group in the form
$$
A=\underbrace{\bigoplus_{i=1}^r\Z}_{\text{torsion-free}}
\oplus
\underbrace{\text{finite}}_{\text{torsion}},
$$
where $r=\ran(A)$ is the rank of the group. This is the
number appearing in the proposition.
\end{prop}

\begin{ejer}
Let
$$
0\to A\to B\to C\to0
$$
be an exact sequence of finitely generated abelian groups.
Then
$$
\text{ran }B=\text{ran }A+\text{ran }C.
$$
See the splitting lemma.
\end{ejer}

Now let us prove the proposition.

\begin{proof}
Take the cellular chain complex
$$
0\to C_k\xrightarrow{d_k}C_{k-1}\xrightarrow{d_{k-1}}
\cdots\xrightarrow{d_2}C_1\xrightarrow{d_1}C_0\to0,
$$
and recall that $Z_n=\ker d_n$, $B_n=\text{img }d_{n+1}$,
and $H_n=Z_n/B_n$. Thus we have the short exact sequences
$$
0\to Z_n\to C_n\xrightarrow{d_n}B_{n-1}\to0,
\qquad
0\to B_n\to Z_n\to H_n\to0.
$$
Using the exercise, we have
$$
\text{ran }C_n=\text{ran }Z_n+\text{ran }B_{n-1},
\qquad
\text{ran }Z_n=\text{ran }B_n+\text{ran }H_n.
$$
Now first by definition and then using the above,
$$
\begin{aligned}
\chi(Y)
&=\sum_{i=0}^{\infty}(-1)^i\text{ran }C_i\\
&=\sum_{i=0}^{\infty}(-1)^i
\big(\text{ran }B_i+\text{ran }H_i
+\text{ran }B_{i-1}\big)\\
&=\sum_{i=0}^{\infty}(-1)^i\text{ran }H_i.
\end{aligned}
$$
Just checking carefully what happens with $i=0$ for
$\ran B_{i-1}$.
\end{proof}

And with this proposition the theorem above follows easily.

\chapter{Homology and homotopy}

\section{The Hurewicz homomorphism}

We begin with some algebraic ideas.

Let $G$ be a group and let $g,h\in G$. The commutator of
these elements is
$$
[g,h]=ghg^{-1}h^{-1}.
$$
Two elements commute if and only if the commutator is trivial.
The \textbf{derived subgroup}, or \textbf{commutator
subgroup}, is
$$
G'=\langle[g,h]\mid g,h\in G\rangle\leq G.
$$
Now note that a homomorphism of $G$ restricts to a
homomorphism of $G'$. That is, the restriction is an
automorphism of $G'$. We say that $G'$ is a
\textbf{characteristic subgroup}. This implies that it is a
normal subgroup, so we take the quotient; the result will be
an abelian group because everything that does not commute has
been identified. This group is
$G^{\text{ab}}:=G/G'$ and is called the \textbf{abelianization}
of $G$.

\begin{obs}
Doing this with a simple group, whose normal subgroups are
itself and the trivial subgroup, the quotient $G/G'$ is the
trivial group. But, in general, the idea is to produce an
abelian group from any group.
\end{obs}

\begin{prop}
Let $G$ be a group and $A$ an abelian group. If
$\varphi:G\to A$ is a homomorphism, then there exists a
unique $\bar\varphi:G^{\text{ab}}\to A$ such that the
corresponding diagram commutes.
\end{prop}

\begin{defn}
The \textbf{Hurewicz homomorphism} is
$$
\begin{aligned}
h:\pi_1(X,x_0)&\to H_1(X)\\
[f]&\mapsto [f],
\end{aligned}
$$
which is well defined because loops are indeed cycles, and it
does not depend on the representative.
\end{defn}

\begin{teo}
If $X$ is path connected, then $h$ is surjective and
$\ker h=(\pi_1(X,x_0))'$. That is, the abelianization of the
fundamental group is isomorphic to the first homology group
of $X$.
\end{teo}

\begin{coro}\leavevmode
\begin{itemize}
\item $\Z\cong\pi_1(S^1,x_0)\cong H_1(S^1)$;
\item $\Z\cong\pi_1(S^1,x_0)\cong H_1(S^1)$;
\item one can also deduce for the torus that
$H_1(T)\cong\Z^2$.
\end{itemize}
\end{coro}

\section{Higher homotopy groups}

Recall that the fundamental group was defined as
$$
\begin{aligned}
\pi_1(X,x_0)
&=\{f:I\to X:f(0)=f(1)=x_0\}
\Big/\simeq\text{ rel }0,1\\
&=\{f:(S^1,s_0)\to(X,x_0)\}
\Big/\simeq\text{ rel }s_0.
\end{aligned}
$$
In this second expression for the fundamental group, we can
think of the group operation as being given by a function
defined on the wedge of two circles, since we must:
(1) traverse $S^1$ at double speed, and then (2) identify the
base point with its antipode.

Thus, for $n\geq2$, we can define
$$
\pi_n(X,x_0)=
\{f:(S^n,s_0)\to(X,x_0)\}
\Big/\simeq\text{ rel }s_0,
$$
where the group operation is defined on the wedge of two
$n$-spheres.

\begin{prop}
For every $n\geq2$ and every pointed topological space
$(X,x_0)$, $\pi_n(X,x_0)$ is abelian.
\end{prop}

\begin{obs}
Each element of $\pi_0(X,x_0)$ represents a path-connected
component of $X$. But $\pi_0(X,x_0)$ has no group structure!
\end{obs}

\begin{prop}
$\pi_n:\Top^*\to\Grp$ is a functor for $n\geq1$. That is, it
is well defined on objects and arrows, sends identities to
identities, and is associative.
\end{prop}

\begin{prop}
There is homotopy invariance here too: if
$\varphi,\psi:(X,x_0)\to(Y,y_0)$ are homotopic, then the
induced maps on homotopy groups are the same. In particular,
if two spaces are homotopic through a base-point-preserving
homotopy, then their homotopy groups are the same.
\end{prop}

The converse is not true. There are spaces with the same
homotopy group which are not homotopy equivalent: $\R^2$ and
$S^2$ have the same fundamental group, but their second
homology groups are different.

The converse is not true. There are spaces with the same
homotopy group which are not homotopy equivalent: $\R^2$ and
$S^2$ have the same fundamental group, but their second
homology groups are different. Exercise \ref{ejer:homol-homot}
gives another example.

Now, if we ask that they have the same homotopy groups and
the same homology groups, will they be homotopy equivalent?

\begin{teo}[Whitehead]
Let $X$ and $Y$ be CW complexes. Suppose that $f:X\to Y$ is
such that
$$
f_*:\pi_n(X,x_0)\to\pi_n(Y,y_0)
$$
is an isomorphism for every $n\geq0$. Then $f$ is a homotopy
equivalence.
\end{teo}

In fact, there are spaces whose homotopy groups are
isomorphic but which are not homotopy equivalent; however, in
these cases the isomorphism is not induced by a function
between the spaces.

\begin{teo}
Let
$$
p:(\tilde X,\tilde x_0)\to(X,x_0)
$$
be a covering between connected spaces. Then
$$
p_*:\pi_n(\tilde X,\tilde x_0)\to\pi_n(X,x_0)
$$
is an isomorphism for $n\geq2$.
\end{teo}

\begin{proof}
Take a loop $f:S^n\to X$ and know that it has a lift. This
proves surjectivity. To see injectivity, take a loop in the
covering whose projection is homotopic to zero. We can lift
the homotopy and obtain that the loop upstairs is also
homotopic to the constant loop.
\end{proof}

\begin{ejem}
The homotopy groups of the torus for $n\geq2$ are trivial by
its universal covering, $\R^2$, which is contractible. In
general, any space which has a contractible covering has
trivial higher homotopy groups.
\end{ejem}

\chapter{Exercises}

\begin{ejer}
Compute the homology groups in the examples from the cellular
homology section.
\end{ejer}

\begin{proof}[Solution]
The original notes contain computations using cellular chain
complexes and Mayer--Vietoris. The pictures have been omitted
in this translated version.
\end{proof}

\begin{ejer}\label{ejer:homol-homot}
Show that $S^1\times S^1$ and $S^1\vee S^1\vee S^2$ have
isomorphic homology groups with integer coefficients but are
not homotopy equivalent.
\end{ejer}

\begin{proof}[Solution]
Their homology groups agree, but their fundamental groups do
not: one is $\Z\times\Z$ and the other is the free group on
two generators.
\end{proof}

\begin{ejer}
Compute the homology groups of the space obtained from
$S^2$ by identifying the north pole and the south pole to a
single point.
\end{ejer}

\begin{proof}[Solution]
The simplest CW decomposition of this space has one 0-cell,
two 1-cells, and two 2-cells. This gives the cellular chain
complex
$$
0\to\Z^2\xrightarrow{d_2}\Z^2\xrightarrow{d_1}\Z\to0.
$$
Since there is only one 0-cell, $d_1=0$. The attaching maps
of both 2-cells send their boundary around both 1-cells, so
$$
\begin{aligned}
d_2:\Z^2&\to\Z^2\\
(1,0)&\mapsto(1,1)\\
(0,1)&\mapsto(1,1).
\end{aligned}
$$
Thus $\img d_2\cong\Z$ and
$\ker d_2\cong\langle(1,0)-(0,1)\rangle\cong\Z$, and hence
$$
H_2(X,\Z)\cong H_1(X,\Z)\cong H_0(X,\Z)\cong\Z,
$$
and $H_n(X,\Z)\cong0$ for $n\geq3$.
\end{proof}

\begin{ejer}
Compute the homology groups of $S^1\times(S^1\vee S^1)$.
\end{ejer}

\begin{proof}[Solution]
The cellular chain complex is
$$
0\to\Z^2\xrightarrow{d_2}\Z^3\xrightarrow{0}\Z\to0.
$$
By an argument analogous to the computation of the homology
of the torus, $d_2=0$, and we conclude that
$H_2(X)=\Z^2$, $H_1(X)=\Z^3$, and $H_0(X)=\Z$.
\end{proof}

\begin{ejer}
Use the Mayer--Vietoris sequence to compute the homology
groups of the space obtained by gluing a M\"obius band $M$ to
the torus $S^1\times S^1$ by a homeomorphism from the
boundary circle of $M$ to the equator $S^1\times\{0\}$ of
the torus.
\end{ejer}

\begin{proof}[Solution]
It is not obvious how to glue them at the cellular level,
because the boundary circle of the M\"obius band is formed by
two edges, which should be glued to a single edge of the
torus. Thus it is technically inconvenient to use the
Mayer--Vietoris sequence in its version for CW complexes.
This is not a problem, because although neither the torus nor
the band is open in our space, they are deformation retracts
of suitable neighbourhoods.

Using the Mayer--Vietoris sequence and the known homology of
the spaces involved, we need the map
$\Phi_1=(i_*,j_*)$ induced by the inclusions
$i:T\cap M\hookrightarrow T$ and $j:T\cap M\hookrightarrow M$.
The intersection of $T$ and $M$ is an equator of the torus and
the boundary circle of the M\"obius band. Thus inclusion sends
a generator of the intersection to one of the generators of
the first homology of the torus, and to twice the generator of
the first homology of the band. In symbols,
$$
\Phi_1(1)=(1,0,2).
$$
Thus $\ker\Phi_1=0$ and $\img\Phi_1\cong\Z$. Therefore
$$
\tilde H_2(X)\approx\Z,
\qquad
\tilde H_1(X)\approx\Z\oplus\Z,
$$
and of course $\tilde H_n(X)\approx0$ for $n=0$ and $n>2$.
\end{proof}

\end{document}
